%% !TEX root = ../manual_gaolv5.tex
% Chapter 14.
%% Chapter 14: Tests and continuous integration. Sources: gaol.tex of f5bf180 lines
%% 778-816 (the tests); doc/tests.md, doc/continuous-integration.md,
%% doc/tools.md, examples/examples.md (section 1), examples/tutorial.md,
%% coverage/README.md, manual/check_examples.py (its header comment).
\chapter{Tests and continuous integration}\label{chap:tests}
\cindex{tests}%
\cindex{continuous integration}%

An interval library is trusted for one property: the bounds it computes
enclose the exact result of each operation, on every machine it is built
for. A test that compares GAOL with itself, or with a floating-point library
running in the same environment, cannot establish it, since the reference
would share the roundings and the rounding direction of what it checks. The
tests of GAOL v5 therefore compare the bounds with results computed apart
from the library and from the floating-point environment, with integers,
with exact rational arithmetic, or with mpmath at 2000 bits of precision, and
they run the library in every state a program may leave the processor in.

In this chapter, we describe how GAOL v5 is tested. The first section gives
the principles the whole suite follows: exact references, tests that fail
without the change they check, the floating-point exceptions, the four
rounding directions, the modes that flush the subnormal numbers, the locales,
and the unit tests of GAOL~4. The second section lists the programs of
\file{tests/} and what each of them checks. The third section presents the
example programs of \file{examples/}, which check what they print, and the
script that compiles the examples of this manual. The fourth section
describes the continuous integration: its platforms, the way it measures the
speed of a pull request, the configurations it refuses or leaves out, and
what these platforms showed. The fifth section lists the scripts of the
repository, and the sixth gives the coverage of the tests. The targets that
build and run the tests in the three builds are described with the
installation (\pxref{sec:tests}).

\section{The principles of the tests}\label{sec:test-principles}
\cindex{tests!principles}%

Built with \option{WITH\_TESTS} (CMake), \loption{with-tests}
(\cmd{configure}) or \code{with-tests} (meson), all off by default, the
programs of \file{tests/} compare the bounds GAOL computes with the exact
results of the operations, independently of GAOL and of the floating-point
environment. The exact results are computed with integers, or were computed
with 2000 bits of precision by mpmath. The tests follow the rounding tests
of Codac. They check, among others, that the basic operations are the
tightest enclosures, that the elementary functions are the tightest or
within one double of them, that every operation gives the same result
whatever the rounding direction the calling code left, that numbers are read
exactly and intervals written rounded outward, that the reverse functions
are accurate in the sense of IEEE 1788-2015, and that each name of
\code{gaol\_ieee1788} is the operation of the standard it names.
\code{test\_ibex\_codac} is another kind of test: it checks GAOL against the
cases of the unit tests that IBEX and Codac run on their own intervals,
which are GAOL's. \cmd{make test} runs the tests, and \cmd{make check} runs
the examples of \file{examples/} too (\pxref{sec:tests}).

\medskip
The suite rests on the following principles.

\begin{itemize}
\item[$\bullet$] \textbf{References computed apart.} A reference value is
  never a result of GAOL, nor of the C library running in the same rounding
  direction. It is an integer, an exact rational, or a value computed by
  mpmath with 2000 bits (400 to 7000 for some tables) and written into a
  header that the test includes: \file{elementary\_values.h},
  \file{reverse\_values.h}, \file{atan2\_rel\_values.h},
  \file{pow\_rel\_values.h}, \file{compoundm1\_values.h} and
  \file{extended\_precision\_values.h} are generated by the scripts
  \file{*\_values.py} of \file{tests/}, run by hand when a table changes
  (\pxref{sec:scripts}). Where a function is correctly rounded by CORE-MATH,
  the test computes its tightest bounds by calling CORE-MATH itself in the
  downward and the upward directions, which it sets apart from GAOL.
\item[$\bullet$] \textbf{Tests that fail without the fix.} Each check
  written with a change of GAOL v5 was run against the sources before the
  change, and the documentation of the tests records how many of its checks
  fail there: 10\,900 of the checks under the modes that flush the subnormal
  numbers, and, with the SSE2 intervals, 45 of the 48 calls that compared the NaN
  bounds of an empty operand, which died in a child process that enabled the
  exception, and 25 of the 45 checks of the operations computed before
  \code{main()}. A check that passes with and without the change proves
  nothing, and the mutation tools of \file{tests/tools/} build GAOL with a
  deliberate error that the tests have to catch (\pxref{sec:scripts}).
\item[$\bullet$] \textbf{The floating-point exceptions.} No operation of GAOL
  may raise the invalid-operation, the division-by-zero or the overflow flag
  for a result without an infinite bound, nor lose a flag or a mask the
  program set (\pxref{sec:fp-exceptions}). With glibc, each operation
  checked is also computed in a child process that enabled the exception
  with \code{feenableexcept()}, which must not die on \code{SIGFPE}: the
  comparison of the NaN bounds of an empty interval once did. Where the
  processor does not trap, or keeps no flag, the test says so and skips
  that part.
\item[$\bullet$] \textbf{The four rounding directions.} The bounds must not
  depend on the rounding direction of the calling code: each operation is
  called with the direction set upward, to nearest, downward and toward
  zero (and on x86 with the x87 and the SSE directions differing, six
  states), and has to give the results it gives when called rounding
  upward, and to leave the direction upward, or as it found it with
  \option{GAOL\_PRESERVE\_ROUNDING} (\pxref{sec:rounding-direction}).
\item[$\bullet$] \textbf{The flush-to-zero modes.} Each mode that flushes the
  subnormal numbers to zero that the processor honours (flush-to-zero and
  denormals-are-zero on x86, FZ and FIZ on ARM, \pxref{sec:ftz}) is set just
  before an operation, alone and with the others: the operation has to give
  what it gives with the modes cleared, and to leave them cleared, or set
  back with \option{GAOL\_PRESERVE\_ROUNDING}. A mode the processor keeps
  without honouring it, as an emulator may, is named and skipped.
\item[$\bullet$] \textbf{The locales.} Numbers have to be read as in the C
  locale, and written back bit for bit, under a locale writing a decimal
  comma (\code{fr\_FR.UTF-8}, \code{de\_DE.UTF-8},
  \code{French\_France.1252}) where one is installed, and the names of the
  functions have to be read in any case of letters under the Turkish
  locales of glibc, where \code{std::tolower('I')} is not \code{i}. The
  reading once never ended under a comma locale, and the test, which ctest
  would otherwise let run with no limit, fails after 5 minutes should it
  hang again.
\item[$\bullet$] \textbf{The unit tests of GAOL~4.} The tests of
  \file{check/}, which checked the results of the operations on chosen
  intervals with CppUnit, are in \file{tests/} and run with the checks of
  \file{unit\_tests.h}: CppUnit is needed no more.\newinvfive{} Each suite
  keeps its class and its tests, a check is named after its test and
  describes a failure by its line, its expression and the values compared,
  and a failure no longer ends its test.
\end{itemize}

\begin{Rem}
Under the Debug C runtime of Visual C++, a failed assertion of the runtime
or of the checked iterators of its library, and an invalid parameter, are
written to the standard error, where the test goes on and fails
(\file{tests/gaol\_tests.h}). The runtime showed them in a dialog box, which
nobody closes on a machine of the continuous integration, and the test hung
without output until ctest stopped it after 300~s. That runtime reports such
an assertion when it writes a subnormal number under denormals-are-zero, then
writes 0: \code{operator<<} clears that mode before it writes, which
\code{numbers} checks there too.\newinvfive
\end{Rem}

\begin{Rem}
Built without exceptions (\loption{disable-exceptions} of \cmd{configure},
\option{-Denable-exception=false} of meson), where an error of GAOL ends the
program, the tests leave out the checks of an exception and of a text the
reader refuses.\newinvfive
\end{Rem}

\medskip
\file{tests/find\_package} builds some of the same tests with an installed
GAOL, and \file{.github/scripts/tests.sh} with a GAOL installed by
\cmd{configure} or meson. \file{tests/fetch\_content} builds them with a GAOL
brought into a project by \code{FetchContent}, and
\file{tests/meson\_subproject} with a GAOL brought in as a meson
subproject,\newinvfive{} which meson refused while \file{meson.build} called \code{add\_global\_arguments()}. Its test \code{flags} checks
that \code{gaol\_dep} gives the compilation the flags of the \code{Cflags} of
\file{gaol.pc}, without which the tests pass all the same, and the link the
options of its \code{Libs} (\option{-mno-daz-ftz}, where the compiler takes
it); its test \code{release} checks that GAOL is built in release in the
default build type of meson, debug: \file{tests/release\_flags.py} checks
that the sources of GAOL and CORE-MATH are compiled with \option{-O3} and
\code{NDEBUG}, and that \code{GAOL\_DEBUGGING} is 0. With \loption{debug},
it checks the other way round that \option{GAOL\_DEBUG} of CMake and
\code{enable-debug} of meson build GAOL and CORE-MATH for debugging, GAOL
built alone and brought into another project.

\section{The programs of the test suite}\label{sec:test-programs}
\cindex{tests!programs}%

Table~\ref{tab:test-programs} lists the programs of \file{tests/} and what
each of them checks. The paragraphs that follow give, for each group, the
cases and the numbers of checks worth knowing, and the platforms where a
check applies.

\begingroup\renewcommand{\arraystretch}{1.1}\small
\begin{longtable}{L{4.4cm}L{9.4cm}}
  \caption{The programs of the test suite}\label{tab:test-programs}\\
  \hline
  \multicolumn{1}{c}{Program} & \multicolumn{1}{c}{What it checks}\\
  \hline\hline
  \endfirsthead
  \multicolumn{2}{l}{\textit{Table~\thetable, continued}}\\*
  \hline
  \multicolumn{1}{c}{Program} & \multicolumn{1}{c}{What it checks}\\
  \hline\hline
  \endhead
  \hline
  \endfoot
  \code{arithmetic} & The basic operations, squares, square roots, integer
    powers, \code{fma} and the cancellative operations are the tightest
    enclosures at every magnitude, from every rounding direction, and the
    $n$-th roots are within 2 doubles of them (4 for an inverse root)\\
  \code{elementary} & The elementary functions at doubles, at intervals and
    at intervals whose images are known exactly; \code{pow}, \code{atan2},
    \code{erf}, \code{erfc}, the integer functions\\
  \code{rounding\_direction} & The same bounds from every rounding direction
    and flush mode; the direction left upward or preserved; the control
    registers; the flags and the masks kept; no \code{SIGFPE} on an empty
    operand\\
  \code{fast\_math\_link} & A program linked with \option{-ffast-math}: the
    flush modes that \file{crtfastmath.o} sets\\
  \code{automatic\_cleanup} & (Linux) The rounding direction set back to
    nearest after \code{main()}\\
  \code{static\_initialization} & About 40 operations computed before
    \code{main()} give their bounds of \code{main()}, bit for bit\\
  \code{test\_ibex\_codac} & The cases of the unit tests of IBEX and Codac,
    against the exact results: 3905 checks\\
  \code{numbers} & Reading and writing: the tightest reading of every text,
    the literals of IEEE 1788-2015, the bases 16, 8 and 2, the five output
    formats, the flags of the streams, the locales, the subnormals under
    the flush modes\\
  \code{other\_functions} & Midpoints, widths, radii, distances, splitting,
    \code{max}, \code{min}, the tools of interval algorithms, the
    comparisons of IEEE 1788-2015, the reverse functions of \code{max},
    \code{min}, \code{sign} and \code{floor}\\
  \code{ieee1788} & \code{gaol\_ieee1788} as a program uses it: \code{pow},
    \code{pown}, \code{rootn}, \code{pownRev}, \code{overlap},
    \code{intervalToText}, \code{mulRevToPair}\\
  \code{core\_math} & The bounds against CORE-MATH called in the rounding
    direction in effect; the \code{fma()} and \code{round()} of the C library, and
    \code{gaol\_ldexp()}; the accurate paths of \code{sin} and
    \code{cos}\\
  \code{compoundm1} & $(1+x)^y-1$ from the bottom up: the 256-bit integers,
    the balls, the fast phase, 5\,854 pairs and 850 boxes\\
  \code{expressions} & The lexer, the parser and the evaluation of the
    strings: every node, every wrong string, the exceptions, the deep
    chains, the allocations that fail\\
  \code{u128} & The two 64-bit halves of the 128-bit integer against the
    native type of the compiler\\
  \code{extended\_precision} & The arguments the x87 unit rounds wrongly:
    fails when GAOL is built on it\\
  \code{reverse} & The relational functions against the reverse functions
    of IEEE 1788-2015: valid and accurate; \code{atan2\_rel\_x},
    \code{atan2\_rel\_y}, \code{pow\_rel}, \code{pow\_exponent\_rel}\\
  \code{debugging} & The headers compiled with \code{GAOL\_DEBUGGING}, and
    \code{GAOL\_DEBUG\_VERBOSE} at the level given to \code{gaol::init()}\\
  \code{version\_file} & (CMake) The reading of \file{VERSION.txt}: byte
    order mark, line ends of Windows, UTF-16, NUL bytes\\
  \code{refused\_finite\_\allowbreak math\_only},\newline
    \code{refused\_fast\_math},\newline
    \code{refused\_positive};\newline
    \code{refused\_options.sh} (autotools);\newline
    \code{refused\_options} (meson) & The compilation refused with
    \option{-ffinite-math-only} and \option{-ffast-math}, accepted
    without; \file{refused\_options.sh} and \code{refused\_options} check as well
    the three cases of the tests \code{refused\_aligned\_new*}\\
  \code{refused\_aligned\_new},\newline
    \code{refused\_aligned\_\allowbreak new\_cxx17},\newline
    \code{refused\_aligned\_\allowbreak new\_positive} (CMake) & The SSE2 intervals on
    32-bit Windows refused without \option{-faligned-new=8}\\
  \code{fp\_strict} & (a project of its own, which the Visual Studio jobs
    configure) Visual C++ and clang-cl refused without \option{/fp:strict}\\
  \code{cpack\_stale\_configure} & (CMake) The warning of CPack when the
    committed \file{configure} was generated for another version\\
  \code{automatic\_configuration} & (CMake) The configuration GAOL decides
    by itself, and the settings of the developers\\
  \code{arithmetic\_operators}, \code{assignment}, \code{constants},
    \code{constructor}, \code{float\_functions},
    \code{generic\_functions}, \code{input\_output},
    \code{interval\_functions}, \code{misc}, \code{non\_arithmetic},
    \code{relations}, \code{reverse\_mappings}, \code{trigonometric} & The
    unit tests of GAOL~4: the results of the operations on chosen
    intervals; \code{misc} checks the names of the namespaces and the
    alignment of the intervals\\
  \code{intervalf}, \code{interval2f} & The intervals of floats, where a
    developer compiles them; skipped otherwise\\
  \code{nodiscard\_\allowbreak discard\_cxx11},\newline
    \code{\dots\_cxx14}, \code{\dots\_cxx17};\newline
    \code{nodiscard\_\allowbreak used\_cxx11},\newline
    \code{\dots\_cxx14}, \code{\dots\_cxx17} & A result of GAOL thrown away
    fails the compilation, in the three standards; a result used compiles\\
  \code{gaol\_performance} & (not a test) The time per operation of GAOL's
    arithmetic and elementary functions\\
\end{longtable}
\endgroup

\subsection{The arithmetic and the elementary functions}

\code{arithmetic} checks, on doubles and intervals of every magnitude,
subnormal doubles and overflows included, that sums, differences, products,
quotients, relational divisions, squares, inverses, \code{abs}, \code{min},
\code{max}, \code{\&} and \code{|} are the tightest enclosures, and so are
square roots and integer powers of either sign:\newinvfive{} of doubles for
$n$ from 2 to 9, 16 and 31 to 33 and for $-1$ to $-4$, $-7$, $-16$ and
$-33$, of intervals for $n$ from 2 to 5, 12 and 31 and for $-1$ to $-3$,
$-6$ and $-17$, $x^n$ beyond the largest double or subnormal included, and
for $n<0$ the limits where $x$ holds 0. The suite once allowed up to $2n$
doubles beyond the tightest bounds for $n>0$, and 8 for $n<0$. Every
interval made of 15 special bounds ($-\infty$, the opposite of the largest double, $-2.5$,
$-1$, the opposites of the smallest normal and of the smallest subnormal
number, $-0$, $+0$, the smallest subnormal and normal numbers, $0.1$, 1, 3,
the largest double, $+\infty$), the empty set and the universe, taken two by
two, have to give with \code{+}, \code{-}, \code{*}, \code{/}, \code{\%},
their compound forms (the interval itself as the operand included), with a
double on either side where an interval is a point, and for \code{sqr},
\code{sqrt} and \code{inverse}, the bounds of a reference of plain doubles:
the least of the corners rounded downward and the greatest rounded upward by
the processor, a zero bound times an infinite one counting as 0. An infinite
or NaN double gives the empty set, no operation gives an unexpected empty
set, and the bounds do not depend on the rounding direction of the calling
code, set to each of the four directions before each operation. Inexact
operations on literals the compiler sees (\code{[0.1] + [0.2]},
\code{[1] / [3]}, \code{sqrt([2])}) have to be the tightest enclosures
whatever the direction: an operation computing to nearest fails them, and so
would a compiler that folded the vectors of the library at compile time,
which none does with the \option{-frounding-math} that \code{gaol::gaol}
gives. The sign of three zero bounds the documentation states, that of
\code{[-1, 0]*[-1, 0]} (SSE2 and NEON intervals only), \code{sqr([-1, 2])}
and \code{[0, 1]/[-2, -1]}, has to be the one of the build: it tells a tie
of the maxima of the products, which no other check did, the others
comparing with \code{==}, which does not see it.

\medskip
The powers are checked further. $x^n$ has to be the tightest over 4\,000
intervals near the ends of the doubles, where the products were once rounded
outward and the SSE2 and the FPU intervals gave other bounds for half of
them; 8\,679 powers that are doubles have to be exact, both bounds (an error
bound computed beforehand made \code{pow([2, 3], 3)}
$\itv{8-2^{-50}}{27+2^{-48}}$). $x^2$ has to be \code{sqr(x)}, bit for bit,
the sign of a zero bound included, in each of nine forms of the integer power
with the exponent 2, over the grid of special bounds and 90\,000 random
intervals, each with the calling code in the four rounding directions and,
for the first 3\,000 of each kind, with each flush mode: 4 of the 9 forms
fail without one call of \code{sqr()}, 5 without another. $(-c)^n$ has to be
$c^n$ or $-c^n$ at 1\,088 points near 1, 1.5 and 3 and 12\,736 more. The
powers for $n$ from $2^{24}+1$ to $2^{32}-1$ have to be the tightest against
mpmath: the lower bound was 1962 doubles below the tightest for
\code{pow([1.0000001], 2\textasciicircum{}32 - 1)}; the same for 46 negative $n$ from
$-2^{31}$ to $-2^{24}$, 44 of which fail without the second-order term of
the correction. The $n$-th roots of $x$ below $2^{-1020}$, subnormal
included, have to be within 2 doubles of the tightest bounds (4 for an
inverse root) for $n$ = 4, 5, 33, 65, 129 and 1025 against exact arithmetic
and up to $2^{32}-1$ against mpmath: \code{nth\_root([2\textasciicircum{}-1074], 33)}
was 1983 doubles below. \code{fma(x, y, z)} has to be the tightest over
random boxes, against \code{std::fma} called downward and upward directly:
GCC had folded the negations GAOL rounds its lower bound with.
\code{cancel\_minus} and \code{cancel\_plus} have to be the tightest, and
\code{y + cancel\_minus(x, y)} to enclose $x$, over intervals of the same
width and one double wider, and at the corners where the differences of the
bounds are within a double of each other or beyond the largest double: GCC at
\option{-O3} computed the terms of their TwoSum after the direction was set
back to upward, where it is no longer exact.\newinvfive

\medskip
\code{elementary} checks \code{exp}, \code{log}, \code{sin}, \code{cos},
\code{tan}, \code{asin}, \code{acos}, \code{atan}, \code{sinh}, \code{cosh},
\code{tanh}, \code{asinh}, \code{acosh}, \code{atanh}, \code{sqrt} and
\code{pow} at doubles, at intervals, and at intervals whose images are known
exactly: extrema, poles, domains, \code{log} of an interval holding no
positive number and \code{atanh} of one with no point of $(-1, 1)$ being
empty, \code{exp(0)}, \code{log(1)} and $1^y$ exact. The functions have to
be the tightest where their value is 0, 1, $\pm\pi/4$, $\pm\pi/2$ or $\pi$,
and \code{cosh}, \code{sinh} and \code{tanh} beyond the largest double and
near 1. \code{pow(x, y)} has to be within one double of the tightest bounds
at points, \code{pow([2], [1023.5])} among them, and over 385 boxes of bases
and exponents in every position about the base 1 and the exponent 0, with
the power itself as lower bound where it is a double: \code{pow([4], 0.5)} is
\code{[2]}. \code{atan2} has to be within one double over 324 boxes about
the axes, the tightest over the boxes with infinite bounds or on an axis,
$\itv{-\pi}{\pi}$ across the half-line $y=0$, $x<0$, and empty at $(0, 0)$.
\code{sin}, \code{cos} and \code{tan} have to be within one double at every
magnitude, \code{sin([1e-10])} and \code{cos([2\textasciicircum{}60])} included, and over
649 intervals next to their extrema and poles; \code{tan([-M\_PI\_2, M\_PI\_2])}
has to be the tightest, $\pm1.63\cdot10^{16}$, its width being the double
below $\pi$ and holding no pole, and \code{tan} within one double over 158
intervals whose bounds are next to two consecutive poles.\newinvfive{}
The integer functions
\code{sign}, \code{trunc}, \code{round\_ties\_to\_even} and
\code{round\_ties\_to\_away} have to be exact whatever the rounding
direction, which \code{std::nearbyint} and \code{std::rint} would not be,
over the halfway values, the whole numbers, their neighbours and the
magnitudes beyond $2^{52}$. \code{erf} and \code{erfc} have to be the
tightest against mpmath at 409 doubles (tiny and subnormal arguments,
changes of binade, the arguments past which $\mathrm{erf}(x)$ is within
$2^{-53}$ of $\pm1$ and $\mathrm{erfc}(x)$ is subnormal, and hard cases of
the \file{erfc.c} of CORE-MATH) and over the intervals between
them.\newinvfive{} The values are in \file{elementary\_values.h}, which
\file{elementary\_values.py} generates.

\subsection{The rounding direction and the floating-point \texorpdfstring{\mbox{environment}}{environment}}

\code{rounding\_direction} calls about 110 operations of GAOL's interface
with the rounding direction upward, to nearest, downward and toward zero
(and on x86 with the x87 and the SSE directions differing): they have to
give the results they give rounding upward, and to leave the direction
upward, or as they found it with \option{GAOL\_PRESERVE\_ROUNDING}; among
them the roots of negative order, the integer powers of large exponents and
\code{modulo\_k\_pi()} next to multiples of $\pi$, which compute after one
check of the direction.\newinvfive{} Products and sums have to be the
tightest in a loop that changes the direction before each of them. The 30
forms of \code{+ - * / \%}, \code{sqr}, \code{sqrt} and \code{inverse}
with intervals, with a double and as compound assignments have to be the tightest enclosures on random operands, in a loop that changes the rounding direction before each of them; on ARM they are run $2^{16}$ times each, in the four directions in turn, right after the program wrote the control register itself.
\code{pow} with a subnormal
result, which CORE-MATH rounds in the direction \code{fegetround()} gives,
has to be the tightest in each state, the x87 unit to nearest and the SSE
instructions upward included: the \code{fegetround()} of glibc reads the x87
unit, and the upper bound was below the exact value.\newinvfive

\medskip
The control registers of the processor are checked apart from the library
(\pxref{sec:rounding-direction}). What GAOL writes to the control word of the
x87 unit and to MXCSR on a 32-bit x86 processor is checked on every platform,
with two variables standing for the registers, in the 64 states of the two
rounding fields and the two flush modes, with other bits that must not change
and exception flags raised meanwhile that must not be lost: only the registers
that differ are written, and each is read once; on such a processor, the
registers themselves are checked in the states it can run. The
SSE2 operations with \option{GAOL\_PRESERVE\_ROUNDING}, which keep MXCSR
alone, are checked in its 16 states: MXCSR is written once when the
direction is not upward or a mode is set, and never read or written
otherwise. On x86-64, \code{get\_rounding()}, \code{set\_rounding()} and
\code{set\_rounding\_and\_flush\_modes()} change neither the exception flags
nor any other bit, and the type through which the library reads and writes
the two registers has to read the words the test reads and write the words
it is given, with four sets of the other bits, so that a read that drops the
flags or a write that forces the masks is seen. On 64-bit ARM, FPCR is read
and written with \code{mrs} and \code{msr}, and its bits other than the
direction and the flush modes (the default-NaN mode, FZ16, the trap-enable
bits, and the AH bit of Armv8.7 where the processor retains it) have to be
the program's after each operation.\newinvfive{} With mingw-w64 on x86-64,
the control word of the x87 unit is read with \code{fnstcw} and is part of
the state compared, since the \code{fegetround()} of mingw-w64 13 and later
joins the fields of the two units and its \code{fesetround()} clears the six
flags of MXCSR. Every operation has to keep the exception flags the program
raised (the five of IEEE 754, and on x86 the six of the SSE control
register) and the masks, from each direction, and after
\code{gaol::cleanup()} and \code{gaol::restore\_rounding()}: with
\option{GAOL\_PRESERVE\_ROUNDING}, the SSE2 operations masked the exceptions
again and cleared the flags.\newinvfive

\medskip
The empty set and the special values must raise nothing
(\pxref{sec:fp-exceptions}). \code{is\_empty()} has to be true for six
empty sets and false for three intervals, and to raise no flag; with glibc,
each is also told empty in a child process that enabled the
invalid-operation exception, which must not die: the comparison of the NaN
bounds with \code{<=} did. The operations of the interface with an empty
operand on either side, about 220 calls in 66 checks, among them one for each
of the 48 calls that compared the NaN bounds and five choices a program makes
with \code{is\_empty()} that GCC for 32-bit ARM and for POWER9 made signaling,
have to give the result of the empty set, raise no flag and not die in the
child process: 45 of these 48 calls died with the SSE2 intervals, and all 48
with the FPU ones.\newinvfive{} So do 18 operations given a NaN double, the products of
intervals with bounds in $\{-\infty, -2, -0, 0, 3, +\infty\}$, and the relational
divisions of intervals with bounds in $\{-\infty, -2, -0, 0, 4, +\infty\}$ (the SSE2 intervals divided $+\infty$
by $+\infty$, or a bound by a zero one, in the half of a register they did
not keep), 13 powers with an exponent of extreme magnitude, 10 powers next to
the largest double, where CORE-MATH's \code{pow} computed $\infty-\infty$,
and 65 loops of 64 relations and 4 loops that count the pairs for which a
relation holds, which a compiler vectorizes or if-converts: GCC 9.4 and 13 at
\option{-O3} with \option{-mfma}, GCC 13 to 15 for 32-bit x86, GCC 14 for
POWER8 and GCC 11 to 15 for 64-bit ARM made some of them raise the
exception with the FPU intervals, and Clang 18 and 20 at
\option{-O3 -march=native} without \option{-frounding-math}, so that the
teeth of these loops depend on the processor the code is compiled for. The
midpoints of five intervals with a bound of $2^{1023}$ or more have to be
right and raise no overflow flag, the sum of the bounds having overflowed;
36 calls of functions that computed beyond the largest double on the way
(\code{sqrt()}, \code{nth\_root()}, \code{atan()}, \code{sin()},
\code{pow()}, the relational functions) are run on 20\,000 draws of operands
each, and a result without an infinite bound has to come with neither the
overflow nor the division-by-zero flag;\newinvfive{} 27 cases such as
\code{sqrt([1, DBL\_MAX])}, \code{tan([DBL\_MAX])} and
\code{pow([2\textasciicircum{}-1074], [1.5 2\textasciicircum{}1014])} have to give their tightest
bounds, computed apart, with neither flag, and the same bounds in a child
process that enabled both exceptions.

\medskip
With a mode that flushes the subnormal numbers set, and the direction upward
or not, products, sums, differences, quotients, squares and \code{exp} with
a subnormal operand or result have to be the tightest, computed apart with
exact rationals (and mpmath for \code{exp(-740)}), and the modes have to be
cleared after the operation, or restored with
\option{GAOL\_PRESERVE\_ROUNDING}: a program linked with \option{-Ofast} gets
them from \file{crtfastmath.o}, and \code{[1e-300] * [1e-20]} was
\code{[0, 0]}.\newinvfive{} Then each of the 96 operations that check the
direction, called with each mode set just before it and each direction, has
to give what it gives with the modes cleared and leave them cleared (with
GCC 13 at \option{-O3}, \code{asinpi} of $\itv{100\cdot2^{-1074}}{100\cdot2^{-1074}}$ was above the
exact value and \code{sqrt([-1e-310, 4])} the empty set), and each of the 51
functions that make no check has to leave the modes as it found them, its
doubles compared by their bits; \code{midpoint()}, \code{split()},
\code{chi()} and the writing of an interval, which make the check, have to
clear them: with the sources before that change, 10\,900 of these checks
fail.\newinvfive{} \code{bisect(ratio)} and \code{is\_bisectable()} have to
give, under each mode and direction, what the bits of the bounds give, told
by integers: 2280 of these checks fail with the SSE2 intervals before the
change, and 4008 with \option{GAOL\_PRESERVE\_ROUNDING}. The AVX-512 path is
taken whole where the machine building has the instructions: its bounds are
those of the SSE2 path bit for bit, but for the sign of a zero bound of
\code{x * d} where the zero bounds of $x$ differ in sign, and it leaves the
direction as the program set it (\pxref{sec:avx512}).\newinvfive{}
\code{gaol::cleanup()} has to set back the direction the first
\code{gaol::init()} found, although a static object set it upward before
\code{main()}, and \code{gaol::restore\_rounding()} to set it to nearest
again as many times as it is called.

\medskip
Three programs complete it. \code{fast\_math\_link} is linked with
\option{-ffast-math} (its link only, with GCC and Clang), which links
\file{crtfastmath.o} and the flush modes it sets, unless
\option{-mno-daz-ftz} keeps it out: where the build gives the link that
option, the modes have to be clear when \code{main()} starts; otherwise the
test sets them, as a plug-in built with \option{-Ofast} does, and products,
a difference and \code{exp} with subnormal operands have to be the tightest
and the modes cleared after them.\newinvfive{} \code{automatic\_cleanup}
(Linux only) leaves the rounding direction downward at the end of
\code{main()}, and GAOL's automatic cleanup has to set it back to nearest,
or to leave it downward with \option{GAOL\_PRESERVE\_ROUNDING}: the
direction is read by a function that a constructor of priority 101 registers
with \code{std::atexit()}, before GAOL initializes itself.
\code{static\_initialization} computes about 40 operations before
\code{main()}, in the initialization of a static object: the constants
(\code{pi()}, \code{one()}, \code{universe()}), products and quotients, the
elementary functions and the reader of strings have to give the intervals
they give in \code{main()}, bit for bit. The static library of the CMake
build is initialized after the program, except with MinGW-w64, and GAOL's
constants were computed by its dynamic initialization: 25 of the 45 checks
failed with the SSE2 intervals, 24 with the others, \code{pi()} giving
\code{[-0, 0]}.\newinvfive

\subsection{Reading and writing}

\code{numbers} checks that \code{textToInterval("0.1")} is the tightest
interval enclosing the number read, and the number itself when it is a
double; that the constants are the tightest enclosures of $\pi$, $2\pi$ and
$\pi/2$; that no \code{const char*} nor \code{nullptr} converts to an
interval, the constructor from one string being gone;\newinvfive{} that the
literals of IEEE 1788-2015 (\code{[entire]}, \code{[ ]}, \code{[1,]},
\code{3.56?1}, hexadecimal numbers) are read as the tightest intervals in
any case of letters, and the numbers in base 16, 8 and 2 (\code{0x10},
\code{0xff.8}, \code{0o17}, \code{0b1.1p3}) as the tightest intervals
enclosing what their digits write, computed from the digits by the test,
beyond $2^{64}$, around $2^{53}$ and beyond the doubles, while \code{0x},
\code{0o8} and \code{0b2} are refused (\pxref{sec:reading-intervals}).
Expressions GAOL cannot read or compute have to throw, and, built with
LeakSanitizer, the parser has to free the nodes of all of them. The
intervals written in decimal in the general, scientific and fixed formats
with 1 to 20 digits have to enclose the intervals, each bound less than one
unit of its last digit away, near the powers of ten too; read back, they
have to enclose the intervals written. The format of the agreeing digits
writes the shared digits only for finite, nonzero bounds of the same sign,
with as many digits before the point and the same exponent
(\code{-1.25\textasciitilde[67, 0]}), and the bounds otherwise. A point
interval is written \code{[a]} when the number is the point itself, and
\code{[l, r]} otherwise, neither bound being the point
(\code{[0.1, 0.1000000000000001]} for \code{interval(0.1)}); the largest
doubles and the subnormals, written with all their digits (309 digits, 1074
after the point), are read back as the point.\newinvfive{} In hexadecimal,
the bounds have to be written in the form of IEEE 1788-2015 (13.4.1) and
read back as the same set, which is the recovery requirement of 13.4.
\code{operator<<} has to leave the precision of the stream as it was,
\code{std::setw} to pad the whole interval, and the texts written under the
flags of the stream and a locale of its own have to be those that \code{operator<<} wrote when it still wrote into a
stream of its own, before the text was made in a character string, except that
the bounds format writes in the C locale and the width
format writes the radius without a sign under \code{showpos}: 33 cases of
four formats. The width and center formats have to write the
\code{midpoint()} and the \code{rad()} of the standard, the radius rounded
upward so that midpoint plus or minus it contains the interval, and never 0
for an interval that is not a point: GAOL wrote $(l+r)/2$ and $(r-l)/2$
rounded to nearest, which do not contain the interval.\newinvfive{} The
empty set has to be \code{[empty]} in the five formats.

\medskip
\code{while (in >> x)} has to stop at the end of the input with
\code{failbit} set and the interval unchanged, and a line that is no
interval has to set \code{failbit} and throw \code{std::invalid\_argument};
blank lines have to be skipped, under \code{std::noskipws} too, and the
stream tied to the input flushed before the first character is read, so that
a prompt shows before the user types.\newinvfive{} Numbers with a million
zeros after their point and an exponent of 7 digits have to be read exactly:
the exponent was cut at 100000. Under a locale writing a decimal comma,
numbers of 5000 to 20000 characters, in decimal, in hexadecimal and in the
uncertain form, have to be read as the tightest enclosures, known without
reading them, and the number of 20000 characters has to take at most 10
times the time it takes under the C locale, and 50~ms more: the reader read
the text again for each of its 125 comparisons or so, in a time quadratic
in its length, and took 50 times as long.\newinvfive{} The cache variable
\option{GAOL\_NUMBERS\_TIMEOUT} gives the test more than its 5 minutes on a
slower build (600~s with GCC and the sanitizers on macOS x86\_64). With
each flush mode the processor honours, the numbers from 0 to the least
normal double have to be read as the tightest intervals enclosing them, and
the doubles among them as themselves: the reader compared each number with
the doubles around it as doubles, and denormals-are-zero reads a subnormal
as 0, so that \code{1e-310} was read as the interval from the greatest
subnormal to the least normal double, which does not enclose
it.\newinvfive{} Under such a mode, \code{<a, b>} of two different
subnormals has to be refused, intervals of subnormals written with 16
digits have to be read back as enclosures once the mode is restored, and
the text has to be the one written without the mode, \code{operator<<}
clearing it before it writes.\newinvfive

\subsection[The functions, the standard and the reverse functions]{The functions, the standard and\\ the reverse functions}

\code{other\_functions} checks the midpoints (of subnormal bounds too), the
widths, the radii (that of the point zero, like its width, being $+0$ whatever
the signs of its bounds), the magnitudes and mignitudes, the Hausdorff distances
(of intervals with infinite bounds too), \code{nb\_fp\_numbers()} across
the two zeros, the splitting, the integer parts, and \code{max(x, y)} and
\code{min(x, y)} of every pair of intervals made of special bounds, with the
bits of the zeros and the subnormals, with the flush modes set or not; with
the SSE2 and NEON intervals, every pair of stored lanes made of special
doubles has to give the bits of the formula of the other kinds, and the
quotients of the NEON intervals with a NaN in a lane of an operand the bits
of \code{interval::emptyset()}. The tools of interval algorithms,
\code{width\_enclosure()}, \code{inflate(r)} and
\code{interval::centered\_form(m, r)}, have to be the tightest enclosures of the
exact width, of $\itv{a-r}{b+r}$ and of $\itv{m-r}{m+r}$, computed with exact
rationals; \code{bisect(ratio)} has to give two parts that cover the
interval exactly, cut strictly inside where \code{is\_bisectable()}, for
ratios from $10^{-9}$ to $1-10^{-9}$, and to raise no overflow exception for
finite bounds next to the largest double, \code{bisect(0.5)} being
\code{split()}, and to throw \code{std::domain\_error} for a ratio outside
$(0, 1)$ or NaN.\newinvfive{} The comparisons of IEEE 1788-2015 have to
give the values of its Tables~10.3 and 10.4, \code{==} and \code{!=} must
not compile on intervals, and each name of \code{gaol\_ieee1788} has to be
the operation of the standard it names, which a wrong translation would not
show at compilation: the eight comparisons over 20\,000 pairs, \code{inf},
\code{sup} and the numeric functions against Table~10.2, the reverse
functions with the arguments in the order of the standard,
\code{mulRev(b, c, x)} being \code{div\_rel(c, b, x)}, and \code{pow}
against the \ieee{pow} of Table~9.1, which GAOL's own \code{pow} is not for
a negative base (\pxref{sec:pow-namespaces}). The reverse functions of
\code{max}, \code{min}, \code{sign} and \code{floor}\newinvfive{} have to be
the hulls of their sets, read from the definitions of the functions apart
from the code, on every interval with bounds among $-\infty$, $-2$, $-1.5$,
\dots, 2, $+\infty$ and the two zeros, on 66 cases written out, and on
20\,000 random intervals of every magnitude, where every sampled double the
definition puts in the set has to be in the result.

\medskip
\code{ieee1788} uses \code{gaol\_ieee1788} as a program does, under
\code{using namespace gaol\_ieee1788;} alone: every name of the standard it
provides is called unqualified, which compiles only if none is ambiguous
with a function of \code{gaol\_core} (Chapter~\ref{chap:ieee1788}).
\code{pow} has to be the standard's with an interval, an \code{int} or a
\code{double} exponent, and \code{pown} and \code{gaol::pow} the integer
power for an exponent of any integer type:\newinvfive{}
\code{pown([2], 3000000000u)} is the overflowing power, which the
conversion of the exponent to a negative \code{int} made a tiny one, and an
exponent beyond the long longs or of 128 bits gives the tightest bounds of
mpmath where the power is a double and within three doubles of them
otherwise. \code{rootn()}, \code{gaol::nth\_root()} and \code{pownRev()}
take an order of any integer type, and \code{pownRev(c, x, p)} for $p=0$ and
$p<0$ has to give the exact roots of the reciprocals of $c$ where they are
doubles, with no solution for the part of $x$ in the gap of $c$.
\code{overlap(a, b)}\newinvfive{} has to be the state the set specification
of Table~10.7 gives, whose quantifiers the test reads over a finite model of
the real line, on 31\,449 pairs: exactly one specification holds for two
nonempty intervals, but for \ieee{meets} and \ieee{metBy} with a point
interval on a bound of the other one, which the standard gives to
\ieee{starts}, \ieee{finishes}, \ieee{finishedBy} or \ieee{startedBy};
\code{disjoint}, \code{equal} and \code{subset} have to be the masks of
states 0xF001, 0x8040 and 0xC3C0 of the standard. The bounds of
\code{gaol\_ieee1788::pow} and \code{gaol::pow} on 117 boxes, each reaching
a branch of the \ieee{pow} of Table~9.1, have to be the bounds the table
gives, checked against the exact power computed with 500 bits by mpmath
(\file{tests/tools/pow/checkrows.py}). \code{intervalToText(x)} has to be a
literal of the standard with the digits of the precision of the intervals,
whatever the global format and the locale, which a grammar of the literals
of Tables~9.5 and 12.2 checks, and \code{textToInterval} has to read it back
as an interval containing $x$: under a comma locale the literal
\code{[a]} wrote \code{[-2,5]}, read as $\itv{-2}{5}$.\newinvfive{}
\code{mulRevToPair(b, c)} has to be the pair of 10.5.5, the part below
first, tight on cases whose parts have exact bounds, and every $x=k/4$ of
the set on every $b$ and $c$ with integer bounds in $[-3, 3]$ has to be in
one of the two parts.\newinvfive

\medskip
\code{reverse} checks the relational functions against the reverse functions
of IEEE 1788-2015 (10.5.4, Table~10.1): \code{sqrt\_rel} (\ieee{sqrRev}),
\code{invabs\_rel} (\ieee{absRev}), \code{nth\_root\_rel} (\ieee{pownRev},
for every integer $p$),\newinvfive{} \code{asin\_rel}, \code{acos\_rel},
\code{atan\_rel}, \code{acosh\_rel}, and \code{div\_rel}, \code{\%} and
\code{\%=} (\ieee{mulRev}). Over the cases of the minimal tests of
libieeep1788 and random ones, each result has to be empty when an argument
is, \emph{valid}, enclosing the hull of the $x$ of $I$ whose image is in
$J$, which the test holds the tightest enclosure of, and \emph{accurate},
within $\ieee{nextOut}(\ieee{tightest}(\ieee{nextOut}(\mathit{arguments})))$
(12.10.1), as the standard recommends for inf-sup types. The bounds of
\code{asin\_rel}, \code{acos\_rel} and \code{atan\_rel} are allowed one
double further, their reduction modulo $\pi$ rounding twice, and those of
\code{nth\_root\_rel} for a negative $p$ two: of the 18\,052 random cases
with $p<0$, 12\,495 have a nonempty result with nonempty accurate bounds,
521 were one double beyond them, 25 two, none more
(\pxref{sec:relational-arithmetic}).
The hulls of \file{reverse\_values.h} are computed from the definition of
each function with exact rational arithmetic, and mpmath for the periodic
ones, not from GAOL nor from the results libieeep1788 expects, which the
generator checks are enclosures of them. \code{atan2\_rel\_y} and
\code{atan2\_rel\_x} (\ieee{atan2Rev1}, \ieee{atan2Rev2}),\newinvfive{}
which libieeep1788 does not test, have a table of their own, generated with
mpmath at 3000 bits: 1\,969 cases, 1\,263 of them with points, whose hulls
are computed from the vertices of the polygon of each convex cone in the
box; each result has to be valid, accurate and within 2 doubles of the
tightest bounds, and about 600\,000 sampled points whose angle is in $c$
have to be in both results. \code{pow\_rel} and \code{pow\_exponent\_rel}
(\ieee{powRev1}, \ieee{powRev2})\newinvfive{} are checked on the 804 cases
of \file{pow\_rev.itl} of ITF1788, random ones and ones whose solutions are
doubles, the hulls computed with mpmath at 2000 bits by cells: the generator
found two results of ITF1788 wider than the tightest enclosure, and by
sampling over 6\,000 random boxes, 54\,305 points whose power is within $c$
have to be kept by both functions.

\medskip
\code{test\_ibex\_codac} checks the cases of the unit tests that IBEX and
Codac run on their intervals, which are GAOL's, against the exact results:
3905 checks.\newinvfive{} A case is a row of a table, written with
\code{gaol::interval} and the functions of GAOL, named by its origin
(\code{IBEX TestArith::add05}, \code{Codac Interval\_operations:176}); the
exact result is computed apart by the scripts of
\file{tests/tools/ibex\_codac} with exact rationals and mpmath at 2000 bits,
and the bounds of GAOL have to enclose it with no double between them and
the tightest bounds (two doubles for \code{atan2\_rel\_x} and
\code{atan2\_rel\_y}). Three sections: the forward tests of IBEX (71 of the
139 tests of \file{TestInterval} and 182 of the 343 of \file{TestArith},
and 3 that IBEX never runs; 1393 checks), its backward contractors, its
generalized division and its distances (197 cases, 1306 checks), and the
tests of Codac (1206 checks; of the 945 checks of its three files, 195 are
already in the tests of GAOL). What only IBEX or Codac have, \code{tan06} of
IBEX, whose expectation is wrong, and the distance to the empty set, whose
convention differs, are left out. What the projects expect is compared with
the exact results by the scripts: 1940 comparisons, no disagreement beyond
their tolerances, and 23 expectations of IBEX with an integer bound that is
not a bound of the exact result (\code{tan02} expects $\itv{-1}{1}$ where it
is $\pm(1+2^{-52})$).

\subsection{CORE-MATH, compoundm1 and the 128-bit integer}

\code{core\_math} checks the bounds of the elementary functions against
CORE-MATH itself, which is correctly rounded in the rounding direction in
effect: the tightest bounds of $f$ at a double $x$ are the values it gives
rounding downward and upward, which the test computes by setting the
direction itself. GAOL's bounds have to enclose them, to be equal above and
within one double below, and to be the tightest where the code gives the
exact value itself. Each function is tried at the ends of its domain, at the
values GAOL treats apart, at the powers of two and their neighbours, at the
subnormals and at random doubles. \code{pow(x, y)} has to be the tightest
at the corner of a box, both bounds, over 84\,070 pairs; where $x^y$ is
rational (8\,981 pairs, 3\,211 of them a double) the reference is $x^y$
itself, computed with integers apart from GAOL and from CORE-MATH, whose
values downward and upward have to be its roundings too, which tells a
fault of CORE-MATH from one of GAOL: CORE-MATH's \code{pow} rounded to
nearest the midpoints of its exact phase on 32-bit ARM, and its square roots
on a 32-bit x86 Windows. \code{exp2}, \code{exp10}, \code{log2} and
\code{log10} have to be the tightest, and the exact values where they are
doubles (\code{exp10} of 0 to 22, \code{log10} of a power of ten up to
$10^{22}$). The functions of Table~10.5 but \code{compoundm1}
have to be the tightest over intervals too, the hull of their image computed
from the values at the bounds of the part of the interval in the domain, at
the points nearest to and farthest from the origin (\code{hypot}), and at
every multiple of $1/2$ within for \code{sinpi}, \code{cospi} and
\code{tanpi}.\newinvfive{} \code{sin} is checked, in the four directions,
against the upstream sources of CORE-MATH compiled unchanged at 19 arguments
of its accurate path, which computes with the 128-bit integer GAOL ports,
and \code{cos} at 45 arguments of its accurate path against its roundings
computed with mpmath at 4\,000 bits, which reach the branches of that path.
First of all, the test checks the \code{fma()} and \code{round()} of the C
library, which CORE-MATH calls: \code{fma()} correctly rounded in the four
directions at six triples, and \code{round()} at
$\pm$\code{0x1.fffffffffffffp-2}, $\pm0.5$, $\pm2.5$ and $2^{52}-1/2$;
the \code{fma()} and \code{round()} of
mingw-w64's own math library on x64, which GAOL refuses, fail 21 and 6 of
these checks under wine.\newinvfive{} The \code{gaol\_ldexp()} CORE-MATH
rounds the subnormal values of \code{erfc} with has to be $x\cdot2^n$
rounded once in the four directions at 20\,000 random pairs. \code{pow} at
the corners where the base is 0 or $+\infty$ or the exponent infinite has to
raise neither the division-by-zero nor the invalid-operation exception, and \code{pow} of a power next to the largest double no
invalid-operation exception either: for $y\log x$ from
709.78263 to 709.78267, CORE-MATH's first phase scaled its approximation by
$2^{1024}$ and computed $\infty-\infty$, at 942 of 1\,001 exponents of the
base 2 from $1024-2^{-13}$ in steps of $2^{-24}$.\newinvfive{} The reading
of a number above the largest double (\code{1e400}, \code{0x1p1024}) has to
be $\itv{\mathit{DBL\_MAX}}{+\infty}$ and raise no overflow flag, which
\code{strtod()} raised, and a number below the least normal double no
underflow flag.\newinvfive{} CORE-MATH, called without GAOL in each of the
six directions of x86, has to round in the direction of MXCSR; on a 32-bit
x86 processor its \code{fegetround()} was the C library's, reading the x87
unit (glibc) or giving $-1$ (Visual C++) when the two differ. \code{erf}
and \code{erfc} have to be the tightest over 80\,000 intervals, and at 22
hard-to-round arguments of the \file{erfc.c} of CORE-MATH.

\medskip
\code{compoundm1} checks $(1+x)^y-1$ of IEEE 1788-2015 (Table~10.5), which
CORE-MATH has not and whose bounds GAOL computes with 64-bit fixed point in
a fast phase, and where that does not settle them with a 128-bit
significand and a bound of the error.\newinvfive{} The test goes from the
bottom up: the count of leading zeros, the 256-bit product against the
schoolbook one, the integer square root, the doubles around an integer; the
sums and products of 20\,000 triples of doubles, whose balls have to hold
the exact value; the balls of $\log(1+x)$ for 493 arguments and of $e^t-1$
for 322, which have to hold the value computed to 256 bits by mpmath and be
no wider than $2^{-118}$ relatively; the fast phase by itself, whose center
has to be within $3/10$ of the bound of its error over 20\,000 pairs from
eight domains (the worst error found is 0.22 of it, in 27 million pairs),
and which answers 74 to 96\,\% of the pairs according to the domain; the
exact division of $y=-1$ at the 633 points of 719 where $1+x$ takes 127 bits
or less. Then the bounds at 5\,854 pairs have to be the greatest double at
most the value and the least double at least it, the values computed by
mpmath with 1\,500 to 7\,000 bits, and on 850 boxes the result has to be
the hull of the values at the four corners and of 0 where the box crosses
$x=0$ or $y=0$, with the limits at $x=-1$ and at infinity, read in the
script and not in GAOL.

\medskip
\code{u128} compiles the two 64-bit halves of \file{gaol/math/gaol\_u128.h},
with which GAOL computes the accurate phases of CORE-MATH where the compiler
has no 128-bit integer (Visual C++, and GCC on 32-bit targets), and compares
every operation with the native type at the ends of the ranges, at the
powers of two, at random values and over every shift count from 0 to 127:
the two have to give the same 128 bits. Where the compiler has no native
type, the test checks the identities the operations satisfy instead.
\code{extended\_precision} checks the results that doubles computed on the
x87 unit round wrongly, which is why \file{gaol/gaol\_config.h} refuses
that unit (\pxref{sec:refused}): \code{exp(-0x1.74910d52d3051p+9)} gave 0
rather than $2^{-1074}$ and \code{tanh(0x1.30fc1931f09c9p+4)} 1 rather
than $1-2^{-53}$, their results being rounded twice. The values are computed
by \file{extended\_precision\_values.py} with mpmath at 5000 bits. Built on
the x87 unit, the test fails on Debian 12 i386 with GCC 12 (22 checks) and
on x86-64 with \option{-mfpmath=387} and GCC 9 (76 of its 300 checks).

\subsection{The expressions}

\code{expressions} goes through every node of the tree of
\file{gaol/parser/gaol\_expression.h} and every way a string can be wrong:
the numbers in every form the lexer takes, the operators and the functions
alone and nested, and the strings the parser has to refuse with an exception
rather than an interval. The intervals have to be read wherever a number may
stand (\code{1+[1,2]}, \code{[cos([0,1]), 2]}), which the two grammars of
GAOL~4 refused, an error anywhere has to stop the reading, and the names
GAOL v5 adds (\code{exp2}, \code{log2}, \code{cbrt}, \code{sign},
\code{trunc}, and every function of GAOL) have to be read in any case of
letters, the lexer taking the longest name.\newinvfive{} Each value is
compared with the same computation written in C++, which the other tests
check against the exact results: what is tested is the lexer, the parser
and the evaluation, not the operations. What GAOL throws has to be exactly
\code{std::invalid\_argument} for what the reader refuses and
\code{std::domain\_error} for a function called with an argument it does
not take, told apart by \code{typeid}, with a \code{what()} that is
\code{file:line: explanation}, the file being the name of the source that
threw with no directory (\pxref{sec:exceptions}).\newinvfive{} Long
expressions have to be evaluated and deleted: chains read from a string of
200000 sums or products and of 20000 differences, and chains of 200000
minus signs and 20000 sines, whose evaluation and deletion made one call
within another per node and overflowed the stack of 8~MB; a leaf at the
bottom of a chain of 20000 nodes has to be visited and deleted less than
64~KiB of stack away from the function that starts the evaluation, where the
recursion took 1.9~MB and 0.96~MB (GCC 15.2, x86-64, Release). A string
nested 9000 levels has to be read, and one nested 100000 levels refused, and
a sum of 10000 terms read with no more than 64 allocations live at once,
where the parser built 20000 nodes before the evaluation.\newinvfive{} A
string holding a NUL character has to be no expression, a null pointer has
to throw \code{std::invalid\_argument} with the names of GAOL and of the
standard, and to make \code{parse\_interval()} return \code{false} without
exceptions, where it died of a segmentation fault. \code{[l, u]} with $l$
above $u$ has to be the empty set, two numbers between the same doubles
being compared exactly, as digits, and the bounds of \code{<a, b>} have to
be doubles: \code{<0.1, 0x1.9999999999999p-4>} is refused.\newinvfive{}
Every allocation of a reading has to be able to fail, each in turn: the
reading has to throw \code{std::bad\_alloc} and free all it allocated,
which the test counts with a replacement of the global \code{operator new},
where flex ended the program with \code{exit(2)}. The checks of several
threads reading at once are commented out, the tests running no thread.

\medskip
\code{debugging} compiles GAOL's headers with \code{GAOL\_DEBUGGING}, which
the Debug builds define, and builds and evaluates an expression whose nodes
call \code{GAOL\_DEBUG\_VERBOSE}, which has to run its command at the level
given to \code{gaol::init()} and not above (\pxref{sec:debugging-facilities}):
\file{gaol/parser/\allowbreak gaol\_expression.h} did not compile with
\code{GAOL\_DEBUGGING}, writing on \code{std::cout} without the
header.\newinvfive

\subsection{The builds and the compilers}

\code{version\_file} (a script of CMake) checks that the reading of
\file{VERSION.txt} ignores a UTF-8 byte order mark, the line ends of
Windows and the blanks around the version, refuses anything but three
numbers, and names a UTF-16 file and a NUL byte in its message, whether
\code{file(READ)} cuts the text at the NUL byte (CMake 3.14.7 and 3.16.3)
or keeps it (3.28.3, 4.4.3); \file{.github/scripts/version-file.sh} runs the
same cases on \cmd{configure}, meson and autoconf in the continuous
integration.\newinvfive{} The compile tests \code{refused\_*} compile a
program including \code{<gaol/gaol>} with \option{-ffinite-math-only} and
with \option{-ffast-math}, put after the flags of interval arithmetic, and
the compilation has to fail with the message of \file{gaol/gaol\_config.h},
which the test looks for in the output of the compiler; compiled with
neither, it has to compile (\pxref{sec:refused}). With
\option{-ffinite-math-only} the empty interval is no longer told empty:
\code{([1, 2] \& [3, 4]).is\_empty()} was false with GCC 9.4 and 13, and
with Clang 18. The CMake build copies and touches the source before each
compilation that has to fail (\file{tests/compile\_error.cmake}), since an
object compiled once stayed up to date. \code{refused\_aligned\_new} and \code{refused\_aligned\_new\_cxx17} compile
\file{gaol/gaol\_config.h} alone with \option{-D\_WIN32 -U\_WIN64}, which stand for
32-bit Windows on any system, without the \option{-faligned-new=8} the flags
carry there, in C++11 and in C++17, and have to fail;
\code{refused\_aligned\_new\_positive} compiles it with the option, and has to
succeed: the C runtime of 32-bit
Windows aligns the memory of \code{new} on 8 bytes, an SSE2 interval needs
16, and a \code{std::vector} of SSE2 intervals stopped the program on its
first \code{movaps}.\newinvfive{} \code{cpack\_stale\_configure} configures
copies of the tree and checks that CMake warns when the committed
\file{configure} says another version than \file{VERSION.txt}, naming the
two versions, and that CPack configures the copy again and names the archive
after the new version. \code{automatic\_configuration} configures GAOL
several times in one directory, with the defaults, \option{GAOL\_SIMD=OFF},
the AVX-512 path refused and asked for, the Debug build type, the settings
of the developers and a cross-compilation, and holds
\file{gaol\_configuration.h} and the messages of CMake to what each asks for
and to what the checks found; it builds nothing and takes about ten seconds.
Only the CMake build has these two scripts: \cmd{configure} and meson are
checked by the audit of the three builds. The six tests \code{nodiscard\_*}
compile a program in each of the three standards with
\option{-Werror=unused-result} and \option{-Werror=attributes}
(\option{/we4834}, \option{/we5030} and \option{/we5051} with Visual C++):
throwing away the result of \code{sqrt(x)} has to fail, and using it has to
compile, since \code{GAOL\_NODISCARD} is \code{[[nodiscard]]} before C++17
with GCC 7 and later and Visual C++ 2019 16.4 and later, and the attribute of
Clang with Clang.\newinvfive

\subsection{The unit tests of GAOL~4}

\code{arithmetic\_operators}, \code{assignment}, \code{constants},
\code{constructor}, \code{float\_functions}, \code{generic\_functions},
\code{input\_output}, \code{interval\_functions}, \code{misc},
\code{non\_arithmetic}, \code{relations}, \code{reverse\_mappings} and
\code{trigonometric} check the results of the operations on chosen
intervals, as Fr\'ed\'eric Goualard wrote them in \file{check/}.
\code{reverse\_mappings} draws its random intervals from a seed of its own
rather than from the process identifier, so that its checks are the same at
each run, and checks the infinite ends of \code{atanh\_rel(J, I)} where $J$
ends at $\pm1$.\newinvfive{} \code{input\_output} checks that the reader
refuses \code{<-inf, -inf>} and \code{<inf, inf>}, whose bounds are not one
value. \code{constructor}, \code{assignment}, \code{arithmetic\_operators}
and \code{relations} also check the integers: \code{interval(n)} and
\code{interval(a, b)} around $2^{53}$, $2^{63}$ and $2^{64}$, for integers of
every width and in the four rounding directions, each operator in both
orders, and \code{set\_contains()} at $2^{53}+1$ and $2^{64}-1$, which the
conversion to a double made wrong.\newinvfive{} \code{relations} checks
\code{overlap()} on each of the 16 states of Table~10.7 and
\code{gaol::lexicographic\_less}, a strict total order that sorts a
\code{std::vector} holding an empty interval, which \code{<} reads past the
end of (\pxref{sec:lexicographic-less}).\newinvfive{} \code{misc} checks the
names of \code{gaol}, which names those of \code{gaol\_core} one by one,
with \code{static\_assert}: the names of GAOL~4 that Codac and IBEX use
compile, and what GAOL uses for itself (\code{get\_rounding()},
\code{gaol\_pown()}, \code{rounding\_state}, \code{the\_null\_expr}) is not
found through \code{using namespace gaol} (\pxref{sec:namespaces}); it
checks that the kernels the operators call are in \code{gaol\_detail} and
not in \code{gaol\_core}, where they made a program ambiguous
(\code{test\_kernels\_are\_\allowbreak not\_in\_gaol\_core}); that the heap
holds a type
as aligned as an interval in \code{std::allocator}, \code{new},
\code{std::vector}, \code{std::list}, \code{std::deque}, \code{std::map} and
\code{std::make\_shared} (\code{test\_heap\_alignment}, which on Linux with a
\code{::operator new} aligning on 8 bytes fails without
\option{-faligned-new=8} with GCC 15 and Clang 21 and passes with it); that
intervals at addresses of 8 modulo 16 compute where the SSE2 intervals keep
their bounds as two doubles (\code{test\_intervals\_\allowbreak
aligned\_on\_8\_bytes});
that a function working on local intervals, entered through a trampoline in
assembly with the stack at each of the four alignments of 4 bytes, returns
with the right bounds on 32-bit Windows with GCC or Clang, where the system
calls its threads with the stack aligned on 4 bytes
(\code{test\_stack\_entered\_\allowbreak aligned\_on\_4\_bytes}, which GCC 14 at
\option{-O2} fails on Linux for 32-bit x86 without \option{-mstackrealign});
and which builds with \option{GAOL\_PRESERVE\_ROUNDING} write MXCSR alone
(\code{test\_rounding\_macros\_\allowbreak write\_mxcsr\_alone}).\newinvfive{}
\code{intervalf} and \code{interval2f} test the intervals of floats where a
developer compiles them (\code{GAOL\_FLOAT\_INTERVALS}), and are skipped
otherwise (exit status 77). \file{check/fpu.cpp}, an empty test,
\file{check/essai.cpp} and \file{check/performances.cpp}, the benchmark of
GAOL~4, are gone, \cmd{make perf} replacing the last.\newinvfive

\subsection{The benchmark}
\findex{gaol\_performance}%

\code{gaol\_performance} (\file{tests/performance.cpp}) measures the time per
operation of GAOL's arithmetic and elementary functions, of the quotient by a
positive double, of the constructor, \code{floor()}, \code{x \&= y},
\code{x | y}, \code{max()}, \code{min()}, \code{x \% y},
\code{x.inverse()}, \code{midpoint()}, \code{nth\_root(x, -2)},
\code{nth\_root(x, 3)} and \code{asinpi()}, of two relations, of the same
operations on doubles, and, in its last rows, of lines of 15 additions, 15
additions and subtractions, 7 products, 7 quotients and 13 operations of the
four kinds between 8 and 16 intervals.\newinvfive{} Its title names the
intervals GAOL computes with, NEON, SSE2 or FPU, and four rows give the
product, the quotient, the square and \code{pow(x, 2)} of operands of random
signs, the other rows having positive operands only, which a processor
predicts: the operands are 1\,024 intervals that the predictor learns by
heart, and the product of random signs of the FPU intervals before their
kernels, which chose a branch by the signs as GAOL~4 did, took 6.4~ns on them
and 12.4 on 16\,384 (Apple M2). The table is Markdown, which the jobs of the
continuous integration copy into their summary. After the table, it prints
the time of the check each operation makes, $1+(2^{-1060}+0)$ not above 1,
and of the same check made of $2^{-60}$: a processor that takes a microcode
assist for a subnormal operand shows it there; with the NEON intervals, the
check is the read of FPCR. It is not a test (\pxref{sec:ci-speed}). The
tools of \file{tests/tools/pow/} and \file{tests/tools/arith/}, which
compare the bounds of two builds of GAOL and build GAOL with deliberate
errors, are not tests either (\pxref{sec:scripts}).

\section{The examples}\label{sec:examples}
\cindex{examples}%

The directory \file{examples/} holds 15 programs that use GAOL v5 the way
interval arithmetic is used in practice, as IBEX, Codac, C-XSC, INTLAB and
the textbooks use it, and the example of GAOL~4, 16 programs in all.
\file{examples/CMakeLists.txt}
is a project of its own, which uses an installed GAOL as any user's project
would:

\begin{cmdshell}
% cmake -S examples -B build-examples -DCMAKE_PREFIX_PATH=@var^prefix~
% cmake --build build-examples
% ctest --test-dir build-examples
\end{cmdshell}

\noindent
The three builds of GAOL build them too, with the GAOL they build, where
they are asked to (\option{WITH\_EXAMPLES}, \loption{with-examples},
\code{with-examples}), and \cmd{make check} (\cmd{ninja check} with meson)
runs them with the unit tests (\pxref{sec:tests}).

\medskip
Examples 01 to 15 check what they print:\newinvfive{} each enclosure is
compared with a value computed apart (mpmath, closed forms, exact
fractions), each proof of uniqueness or of emptiness is tested, and each
claim of the output (``4 times too wide'', ``the excess is divided by 100'')
is proved with interval arithmetic. A failed check prints
\code{FAILED: <claim>} and the program returns a failure, so that ctest runs
the examples as tests. The checks test containments and qualitative facts,
never digits: they pass on every build. Example 16 checks its output too,
with the digits of an enclosure that no operation rounds. All 16 examples
compile without a warning (\option{-Wall -Wextra}, GCC 9.4 and Clang 18)
and pass their checks with the four builds tried: the SSE2 intervals, the
FPU intervals (\option{GAOL\_SIMD=OFF}), \option{GAOL\_PRESERVE\_ROUNDING=ON},
and the library compiled by Clang 18. Each runs in less than 0.3~s with the
default build (06 takes 1.3~s with \option{GAOL\_PRESERVE\_ROUNDING}). Their
outputs are the same on every build, apart from the sign of a zero bound
(\code{[-0, 4]} with SSE2, \code{[0, 4]} with the FPU code), from the report
of the rounding policy in 13, and from a digit or two where a program
computes its own doubles (the preconditioner of 10). Every program follows
the rules of GAOL v5: \code{\#include <gaol/gaol.h>}, no call to
\code{gaol::init()}, and \code{gaol::cleanup()} right after the last use of
GAOL. Three headers are shared: \file{box.h}, a vector of intervals with
the names IBEX and Codac use for boxes (\code{widest()},
\code{max\_width()}, \code{mid()}, \code{volume()},
\code{bisect(i, ratio)}, \code{\&}, \code{|}), since GAOL has no box type
and every algorithm on several variables needs one; \file{dual.h}, a
forward-mode automatic differentiation over intervals, with which a function
written once as a template gives its natural extension with \code{interval}
and its gradient enclosure with \code{Dual<N>}; and \file{affine.h}, the
affine forms of example 12. Table~\ref{tab:examples} lists the programs.

\begingroup\renewcommand{\arraystretch}{1.1}\small
\begin{longtable}{L{4.3cm}L{9.5cm}}
  \caption{The examples of \file{examples/}}\label{tab:examples}\\
  \hline
  \multicolumn{1}{c}{Program} & \multicolumn{1}{c}{What it computes and shows}\\
  \hline\hline
  \endfirsthead
  \multicolumn{2}{l}{\textit{Table~\thetable, continued}}\\*
  \hline
  \multicolumn{1}{c}{Program} & \multicolumn{1}{c}{What it computes and shows}\\
  \hline\hline
  \endhead
  \hline
  \endfoot
  \code{01\_first\_steps} & Intervals built every way, their numbers, sets
    and relations: containment, \code{interval(0.1)} against
    \code{textToInterval("0.1")}, \code{mid()} against \code{midpoint()},
    no \code{==}, the certainly relations\\
  \code{02\_decimals\_and\_rump} & 0.1 summed ten times, $\pi$ by Machin,
    Rump's polynomial: decimals are not doubles, \code{sqrt(2.0)} is the C
    library's, the width as an alarm\\
  \code{03\_dependency\_problem} & $x-x$, $x(1-x)$ three ways, $(x-1)^5$,
    subdivision, Goldstein-Price, rotations: the dependency problem,
    single-use forms, \code{sqr} against \code{x*x}, the linear convergence
    of subdivision, the wrapping effect\\
  \code{04\_centered\_form} & $x\cos x$ and Chebyshev $T_5$, natural and
    mean-value forms: automatic differentiation over intervals, the
    quadratic convergence of the centered form, the monotonicity test\\
  \code{05\_interval\_newton} & A root of $\sqrt{x}+(x+1)\cos x$, the 7 roots
    of $\sin x - x/10$, $\pm\sqrt2$: interval Newton, uniqueness proofs,
    the extended division, the trap of $f(m)/F'(X)$ when $0\in F'(X)$\\
  \code{06\_global\_optimization} & The global minima of Goldstein-Price and
    of the six-hump camel: Moore-Skelboe branch and bound, a priority queue
    with an explicit comparator, the cluster effect, mean-value form and
    monotonicity test\\
  \code{07\_contractors} & IBEX's backward arithmetic, $x^2+y^2=1$ with
    $y=e^x-1.5$: forward-backward contraction (HC4Revise) with GAOL's
    reverse functions, their names in GAOL, IEEE 1788 and IBEX\\
  \code{08\_sivia} & Codac's SIVIA example, IBEX lab2 and lab3,
    $\sqrt x+\sqrt y\in\itv{0}{1}$: set inversion, guaranteed area brackets,
    contractors in pavers, the domain trap of a library without
    decorations\\
  \code{09\_parameter\_estimation} & IBEX lab5, range-only localization from
    three beacons: bounded-error estimation by set inversion, measurements
    read as text, a contractor with \code{sqrt\_rel}\\
  \code{10\_krawczyk} & Two circles of IBEX's solver, the roots of
    $\lambda+e^{-\lambda}$: a branch-and-prune solver with the Krawczyk
    operator, why IBEX bisects at 45\,\%, a false proof from a wrong
    Jacobian\\
  \code{11\_ode\_enclosure} & $x'=-x^2$, Codac's Lohner system,
    $\int\exp(-x^2)$: Picard's a priori enclosure, Taylor steps with a
    remainder, the mean-value form against the wrapping of $x+hf(x)$,
    interval Riemann sums\\
  \code{12\_affine\_arithmetic} & $x-x$, $x(1-x)$, $(x-1)^5$, $(x-y)^2$,
    rotations with affine forms: a sound affine arithmetic on GAOL alone,
    when it beats intervals and when it loses to them\\
  \code{13\_rounding\_environment} & The program's own doubles,
    \code{printf}, \code{lrint}, \code{strtod}, threads, \code{cleanup()}:
    what GAOL's rounding direction does to the rest of the program, and how
    to live with it (\pxref{sec:upward-effects})\\
  \code{14\_generic\_programming} & One template evaluated with
    \code{double}, \code{interval} and \code{Dual<1>}: argument-dependent
    lookup, \code{using gaol::pow}, \code{T(0.0)}, the standard algorithms
    and containers with intervals\\
  \code{15\_text\_and\_ieee1788} & Intervals read and written as text, one
    contractor step with the names of IEEE 1788: the literals GAOL reads,
    reading a file, exact output, the namespace \code{gaol\_ieee1788}\\
  \code{16\_Goldstein\_Price} & A million evaluations of the Goldstein-Price
    function, timed: the example of GAOL~4, with the true function; the
    enclosure contains the range, 231 times wider\\
\end{longtable}
\endgroup

\begin{Rem}
The example of GAOL~4 dropped the $+1$ of the Goldstein-Price function,
$(x+y)^2$ instead of $(x+y+1)^2$, and so did the manual of GAOL~4, which
showed the same program and called the enclosure the range of $f$, whose
true range, about $\itv{3}{635255}$, is about 237 times narrower than the
enclosure $\itv{-56254330}{94177270}$. Both now compute the true
function.\newinvfive{} Its enclosure, $\itv{-87881320}{147125080}$ over
$\itv{-2}{2}^2$, contains the range $[3, 1015690.27\ldots]$ and is 231 times
wider than it, for lack of a link between the occurrences of $x$ and $y$: the
dependency problem (\pxref{sec:dependency}).
\end{Rem}

\medskip
A short tutorial, \file{examples/tutorial.md}, writes four of the algorithms
of the examples in programs of a few dozen lines, the condensed forms of
examples 03, 05, 07 and 06: the range of a function by subdivision, the
interval Newton method, a contractor, and a branch and bound. Each program
is compiled with the flags GAOL gives (\pxref{sec:compiling}):
\begin{cmdshell}
% c++ -std=c++11 -O2 $(pkg-config --cflags gaol) prog.cpp -o prog \
      $(pkg-config --libs gaol)
\end{cmdshell}
Its outputs are the same with the SSE2, the NEON and the FPU intervals.
The first program evaluates $x(1-x)$ over
$\itv{0.4}{0.6}$, whose range is $\itv{0.24}{0.25}$: the formula written as
it is gives an enclosure 20 times too wide, $1/4-(x-1/2)^2$, where $x$
occurs once, gives the range itself, and cutting $x$ into $2^k$ pieces with
\code{bisect(0.5)}, whose images are joined by the hull \code{|}, narrows
the enclosure linearly. The second encloses $\pm\sqrt2$ by the interval
Newton method, with \code{div\_rel} on the negative and the positive parts
of $f'(X)$, which keeps the gap around the midpoint that the division
\code{/} would fill. The fourth finds the minimum of $x^4-4x^2+x$ over
$\itv{-3}{3}$ by a branch and bound with a priority queue on the lower
bounds, in 19933 cuts, and prints
$\itv{-5.44419306660636}{-5.444192066610894}$, which contains the minimum
$-5.4441920666108974\ldots$ computed by mpmath; the natural form converges
slowly, and example 06 shows that on the six-hump camel function it needs
25\,057 bisections to bring the enclosure of the minimum within $10^{-3}$,
where the mean-value form with the derivative of \file{dual.h} reaches
$10^{-9}$ in 1\,127, and 261 with a monotonicity test. The third program is
the contractor of Example~\ref{ex:contractor}.

\begin{Ex}\label{ex:contractor}
The forward-backward contractor of $x^2+y^2=1$ evaluates the expression tree
forward, from $x$ and $y$ to the sum, which it intersects with 1; it then
goes backward, and each node is intersected with what the others allow. The
relational functions of GAOL are the backward operations:
\code{sqrt\_rel(a, x)} is the hull of the $x$ of $X$ with $x^2$ in $a$, both
signs included, where \code{sqrt(a)} would lose the solutions with $x<0$
(\pxref{sec:relational-arithmetic}).
\begin{example}
#include <iostream>
#include <gaol/gaol.h>
using namespace gaol;

// The contractor of x*x + y*y = 1: forward-backward evaluation of the tree
// s = a + c, a = sqr(x), c = sqr(y), with s = 1. Returns false when the box
// holds no solution.
bool contract(interval& x, interval& y)
{
  // forward: an enclosure of each node over the box
  interval a = sqr(x);
  interval c = sqr(y);
  interval s = (a + c) & interval(1.0);     // the root must be 1
  if (s.is_empty()) {
    return false;
  }
  // backward: s = a + c gives a = s - c and c = s - a ...
  a &= s - c;
  c &= s - a;
  // ... and a = sqr(x) gives x = +-sqrt(a): sqrt_rel() keeps both signs, and
  // intersects with x itself, where sqrt(a) would keep the positive one only
  x = sqrt_rel(a, x);
  y = sqrt_rel(c, y);
  return !x.is_empty() && !y.is_empty();
}

int main()
{
  interval x(0.8, 3.0), y(-3.0, 3.0);
  std::cout << "box           " << x << ' ' << y << '\n';
  contract(x, y);
  std::cout << "contracted    " << x << ' ' << y << '\n';
  interval u(2.0, 3.0), v(2.0, 3.0);
  std::cout << "box           " << u << ' ' << v << '\n'
            << (contract(u, v) ? "may hold a solution\n"
                                : "proved: no solution\n");
  gaol::cleanup();
  return 0;
}
\end{example}
The program prints the following lines.
\begin{onscreen}
box           [0.8, 3] [-3, 3]
contracted    [0.8, 1] [-0.6000000000000001, 0.6000000000000001]
box           [2, 3] [2, 3]
proved: no solution
\end{onscreen}
The first box is cut to $x$ in $\itv{0.8}{1}$ and $y$ in
$\itv{-0.6}{0.6}$, the bounds being enclosures: the points of the box off
this part are not on the circle. The second box is proved empty of
solutions, its \code{s} being empty. In \code{gaol\_ieee1788},
\code{sqrt\_rel(c, x)} is \ieee{sqrRev(c, x)}, \code{div\_rel(c, b, x)} is
\ieee{mulRev(b, c, x)}, with its arguments in another order, and
\code{nth\_root\_rel(c, n, x)} is \ieee{pownRev(c, x, n)}.
\end{Ex}

\medskip
\findex{check\_examples.py}%
The examples of this manual are checked as well.
\file{manual/check\_examples.py} compiles each example of
the files of \file{manual/v5/chapters/} that shows what it
prints, in its lines \code{@outputs}, with an installed GAOL and the flags
of \cmd{pkg-config gaol}, runs it, and compares what it prints with the
manual, line by line:\newinvfive{} the examples once printed 1 and 0 where
the manual showed \code{true} and \code{false}, and \code{chi()} of the
empty set printed \code{-nan} with GCC where it showed \code{nan}. An
example holding a \code{main()} is compiled as it is; the others are put in
a \code{main()} of their own, after \file{<gaol/gaol>}, \file{<gaol/parser/gaol\_expression.h>} and six headers of the standard library (\file{<cmath>}, \file{<iomanip>},
\file{<iostream>}, \file{<limits>}, \file{<sstream>}, \file{<string>}), with \code{using namespace std} and
\code{using namespace gaol}, or \code{gaol\_ieee1788} in the chapter of its
names. A line printed that is too long for the page may be shown on several
lines of output, each break standing for a blank. Each example is a program
of its own: a format or a precision one sets does not reach the next. A
compiler error points to the line of the example in its file of \file{chapters/}, and
every output line of the manual has to be in an example the script reads;
what the manual shows otherwise, in an \code{onscreen} block after a
program or in a comment, is not checked. The compiler is \env{CXX}, or
\cmd{c++}, and \env{PKG\_CONFIG\_PATH} has to lead to the \file{gaol.pc} of
the GAOL to check; the exit status is 1 when an example does not compile,
fails, or prints something else than the manual shows. It is no test of the
builds: the job Ubuntu 24.04 x86\_64 GCC of the continuous integration runs
it, in C++11, on the GAOL it installs, so that a change of the library that
changes one of the outputs shows.

\section{The continuous integration}\label{sec:ci}

The workflows of \file{.github/workflows/} build GAOL with CMake and run the
tests, with \cmd{make test} (\cmd{cmake -{}-build <build> -{}-target test},
\cmd{RUN\_TESTS} with Visual Studio). They run for each pull request, and
for each push to the branches \code{master}, \code{MATH-CORE} and
\code{configure-clean} only: a push to another branch runs nothing until it
is proposed, which halves the jobs of a pull request;\newinvfive{}
\code{workflow\_dispatch} runs them on any branch. \file{performance.yml} and
\file{manual.yml} are the exceptions: the first runs for the pull requests
that change what it measures only (\pxref{sec:ci-speed}), the second for the
pushes and pull requests that change \file{manual/} or itself. The badges of \file{README.md} show
their state on \code{configure-clean}. Table~\ref{tab:ci-platforms} gives
the systems, the processors and the compilers they build on; the list of
the platforms GAOL supports is in the chapter on the installation
(\pxref{sec:platforms}).

\begin{table}[htbp]
  \centering\renewcommand{\arraystretch}{1.1}\small
  \caption{The platforms of the continuous integration} \label{tab:ci-platforms}
  \begin{tabular}{L{3.2cm}L{3.1cm}L{7.3cm}}
    \hline
    \multicolumn{1}{c}{System} & \multicolumn{1}{c}{Processors} &
    \multicolumn{1}{c}{Compilers and builds}\\
    \hline\hline
    Ubuntu 24.04, 26.04 & x86\_64, arm64 & GCC, Clang; the sanitizers;
      \option{-Wall -Wextra -Werror}; CMake 3.14; GCC 11 and Clang 14 of
      Ubuntu 22.04; GCC 9; a shared library; autotools, meson\\
    Debian 12, 13 (containers) & amd64, arm64, armhf; Debian 12 on i386 &
      GCC; built for size on Debian 13 armhf; autotools and meson on i386
      and armhf\\
    manylinux\_2\_28 & x86\_64, aarch64 & CMake\\
    Alpine (musl) & x86\_64, aarch64 & CMake\\
    Debian 13 under qemu & s390x, ppc64le, riscv64 & CMake; Ubuntu 24.04
      under qemu for POWER9 (GCC 13, \option{-mcpu=power9}), also built
      for size\\
    macOS 15, 26, 27 & arm64, x86\_64 (natively or under Rosetta), a
      universal build & AppleClang, the Clang of LLVM, GCC; CMake,
      autotools, meson, the generator of Xcode; the sanitizers\\
    Windows & x86, x64, arm64 & Visual Studio 2022 and 2026 (Release and
      Debug); clang-cl; MinGW-w64 15 and 14 (x86, x64), 13 and 12 (x86);
      MSYS2 UCRT64 (GCC), CLANG64 and CLANGARM64 (Clang), MINGW32
      (autotools, meson); Cygwin (GCC, x86-64); Clang with the target of
      Visual C++ and the driver of GCC; meson with Visual Studio\\
    \hline
  \end{tabular}

\end{table}

\paragraph{Linux.} The jobs of Ubuntu build with GCC and Clang, also with
the address and undefined behaviour sanitizers, with
\option{-Wall -Wextra -Werror}, with CMake 3.14, with the self-test of the
check of the assembly (\file{tests/tools/asm-order/check\_order.py}
\option{-{}-selftest}, in Python alone), and the check of the library itself
with GCC 13.3 and Clang 18, with and without the rounding direction
preserved,\newinvfive{} with the compilers of Ubuntu
22.04 (GCC 11 and Clang 14 on x86\_64, GCC 11 on arm64), whose images GitHub
retires by April 2027, with GCC 9, which has no
\code{\_\_builtin\_roundeven()}, and as a shared library
(\option{BUILD\_SHARED\_LIBS}, \loption{enable-shared},
\option{-Ddefault\_library=shared}: the static library is the default of the
three builds).\newinvfive{} The test \code{cpack\_stale\_configure} also
runs with \code{CC="ccache gcc"} and a Ninja outside \env{PATH}, which it has
to pass on to the copies of the tree it configures, and the tests run with
the \code{FE\_*} values of Cygwin (newlib, 0 to 3) put over those of glibc
for the sources of CORE-MATH. The containers build Debian 12 and 13 on
amd64, arm64 and armhf, Debian 12 on i386, manylinux\_2\_28 and Alpine
(musl) on x86\_64 and aarch64, Debian 13 under qemu on s390x, ppc64le and
riscv64, and Ubuntu 24.04 (GCC 13) under qemu for POWER9; two jobs build for
size (\code{MinSizeRel}, \option{-Os}) on Debian 13 armhf and for POWER9,
where GCC calls the intersection rather than inlining it, and
\code{rounding\_direction} checks the choices a program makes on
\code{is\_empty()} of an empty interval there too; one job on Debian 13
amd64 tests the empty sets before the bounds in the relations
(\option{-DGAOL\_RELATION\_BOUNDS\_FIRST=0}, the order of RISC-V), where GCC
14 raised the invalid-operation flag for an empty operand.

\paragraph{macOS.} The jobs build macOS 15 and 26 on arm64 and x86\_64,
natively or under Rosetta, macOS 27 on arm64 (on the runner labelled
\code{xcode-27}; there is no job of macOS 27 on x86\_64 under Rosetta, GitHub
having no runner for it), and a universal build (arm64 and x86\_64: the
macros of each part, the tests of the arm64 part natively and the test
programs of the x86\_64 part under Rosetta, and the intervals each part says
it has)\newinvfive{} with CMake, with autotools (\option{-arch} twice in
\env{CFLAGS}, \env{CXXFLAGS} and \env{LDFLAGS}) and with meson, on macOS 15
and 26; the generator of Xcode, with the list \code{arm64;x86\_64} and
without any, where CMake gives Xcode the architecture of the machine; and
the sanitizers with AppleClang, LLVM's Clang and GCC.

\paragraph{Windows.} Visual Studio 2022 and 2026 build on x86, x64 and
arm64, in Release and Debug. On x64 and x86 the intervals are the SSE2 ones
(on x86 their two bounds are two doubles, which need no aligned memory), on
arm64 the NEON ones, with Visual C++ and clang-cl, which a step of each job
checks.\newinvfive{} A job on x86 and two on arm64 ask for the FPU intervals
(\option{GAOL\_SIMD=OFF}), so that their code, and the rounding direction
they keep through the x87 unit and MXCSR, stay covered. Three x64 jobs ask for
the AVX-512 path (\option{GAOL\_PREFER\_AVX512}), two with Visual C++ and one with clang-cl, which they take where the
runner has the instructions, with a warning where it has not; since the
runners have them or not according to the machine GitHub gives, five more
jobs run the unit tests that compute with intervals under Intel SDE, which
emulates a processor with them
(Visual Studio 2022 and 2026, Release and Debug, and clang-cl). clang-cl,
the Clang of Visual Studio, builds x64 in Release and Debug with Visual
Studio 2022, and x86, x64 and arm64 with Visual Studio 2026; on x64 it
defines \code{\_\_x86\_64\_\_} and \code{\_WIN32} but not \code{\_\_WIN32\_\_},
which upstream CORE-MATH tests to read the rounding direction, and on x86
GAOL gives it \option{-msse2}, without which it computes doubles on the x87
unit. MinGW-w64 15 builds x86 and x64, Release and Debug, with the SSE2
intervals on both and \option{-faligned-new=8} on x86, which a step checks;
14 builds x86 and x64, 12 and 13 x86; on x86, GCC 15.2 builds with
\option{GAOL\_PRESERVE\_ROUNDING} and \option{-mfma}, and GCC 13.2 with
\option{GAOL\_FMA=OFF}, the \code{fma()} of mingw-w64 11 on the x87 unit.
MSYS2 builds UCRT64 (GCC, and GCC again with
\option{GAOL\_PRESERVE\_ROUNDING}, with and without \option{GAOL\_FMA}: with
\option{-mfma}, the rounding direction goes through MXCSR alone, which
Linux, clang-cl and Visual C++ x64 with \option{/arch:AVX2} run), CLANG64
(Clang, and Clang again with the direction preserved, which takes the
AVX-512 path where the runner has the instructions, GCC for Windows never
taking it) and CLANGARM64 on the arm64 runners (the NEON intervals). Cygwin
builds GCC for x86-64, Release and Debug, with
\option{-mfma -mprefer-vector-width=128}, all the tests passing but those skipped there (\code{intervalf} and
\code{interval2f}, which no build compiles, and \code{cpack\_stale\_configure},
the CPack of Cygwin making no archive), and a configuration with
\option{GAOL\_FMA=OFF} that \file{gaol/gaol\_config.h} has to
refuse.\newinvfive{} Clang with the target of Visual C++ and the driver of
GCC (\cmd{clang++}, not clang-cl) builds in the environment of Visual Studio
with CMake and meson for x64 and x86, and with autotools for x64, the
library in C++14 there, the standard library of Visual C++ needing it. With
GCC on x64 and x86 in Release, a step reads the objects with \cmd{objdump}
(\file{.github/scripts/check-gcc-windows-fma.sh}):
\file{gaol\_interval\_core.cpp} and the \code{cos} of CORE-MATH
have to hold fused multiply-add
instructions, and no object an instruction on a \code{\%ymm} or \code{\%zmm} register on x64, which GCC would spill to a stack Windows x64 aligns on 16
bytes (GCC bug 54412); on x86 such lines are listed and noticed, not an error.

\paragraph{The three builds.} The workflows also build GAOL with autotools
and meson on Ubuntu (x86\_64, arm64), Debian (i386, armhf), macOS (arm64,
x86\_64), MSYS2 (UCRT64 for x64, MINGW32 for x86) and Cygwin, with meson and
Visual Studio 2022 (x64, x86) and 2026 (x64, arm64), and with meson 0.53.2, the oldest
the meson build takes; run the tests, and build them again with the GAOL
they install, which is the only library installed, CORE-MATH being compiled
into it; build GAOL as a part of another project brought in by
\code{FetchContent} and as a subproject of a meson project, whose code
compiled with \code{gaol\_dep} has to get the flags of \file{gaol.pc}, and
check that GAOL and CORE-MATH are compiled as in Release where the project
gives no build type, and that \option{GAOL\_DEBUG} and \code{enable-debug}
build them for debugging;\newinvfive{} check that the three builds agree on
Ubuntu (x86\_64 and arm64, GCC and Clang), macOS (arm64, x86\_64), Debian (i386,
armhf) and MSYS2 (UCRT64, MINGW32) (the audit of \file{.github/audit},
\exref{sec:scripts}); check that
\cmd{configure}, meson and autoconf read a \file{VERSION.txt} that starts
with a byte order mark or has the line ends of Windows, and that meson reads
it with no Python on \env{PATH}; check that \cmd{make distclean} gives the
source tree back as git has it after the autotools have built, installed
and tested GAOL in it; and check the refusals of \exref{sec:ci-refused}.
Jobs of each build restore the rounding direction
(\option{GAOL\_PRESERVE\_ROUNDING}): Ubuntu x86\_64 GCC and Clang (which
take the AVX-512 path where the runner has it) and arm64 GCC and Clang,
Debian i386 and armhf, macOS arm64, Visual Studio x86, x64 and arm64, clang-cl x64, MinGW-w64 x86, MSYS2 UCRT64, CLANG64 and CLANGARM64, autotools and
meson. Jobs of autotools and of meson on Ubuntu x86\_64 build GAOL without
exceptions and run the examples and the tests, which leave out there the
checks of an exception. The jobs built in Release print the time per
operation in their summary, a table that no other run compares.

\paragraph{Which intervals.} On 64-bit ARM processors the intervals are the
NEON ones by default (\option{GAOL\_SIMD}; \option{GAOL\_ARM64},
\loption{enable-arm64} and \option{-Denable-arm64} were an option of their
own),\newinvfive{} so that the jobs of Ubuntu arm64, macOS arm64, the
aarch64 containers, and those of autotools and meson on arm64 test them
(\pxref{sec:neon}). A few jobs build the FPU intervals instead, so that
their code stays covered, the preserved rounding direction included: the
GCC and the Clang of Ubuntu 24.04 arm64, both with the direction preserved
too, macOS 15 arm64, and autotools and meson on Ubuntu 24.04 arm64. The GCC
and the Clang of Ubuntu 24.04 x86-64 build them too, GCC without and with
the direction preserved and Clang with it: the compound products and
quotients of the FPU intervals hold the direction only by the order of their
stores, which the checks of \code{rounding\_direction} with the direction
preserved make a rule.\newinvfive{} A build that fell back to the FPU
intervals without saying so would keep every test green, so the jobs check
which intervals the configuration took, in a step ``The intervals''
(\file{.github/scripts/check-intervals.sh}, which reads
\code{GAOL\_USING\_\allowbreak NEON\_INSTRUCTIONS} and
\code{GAOL\_USING\_\allowbreak SSE2\_INSTRUCTIONS} in
\file{gaol\_configuration.h} and
fails if they are not the expected ones): the NEON intervals on 64-bit ARM,
the SSE2 intervals on x86-64 and on 32-bit x86 (Visual C++, clang-cl, MinGW-w64, MSYS2 MINGW32, Clang for Visual C++ and Debian 12 i386), and
the FPU intervals where the option is off and on armhf, s390x, ppc64le and
riscv64; two jobs give the SSE2 intervals of Linux the form they have on
x86 with Visual C++ (\option{-DGAOL\_SSE2\_UNALIGNED\_STORAGE=ON}).

\paragraph{Warnings and headers.} Five jobs build the library, the sources
of CORE-MATH included, the tests and the examples with
\option{-Wall -Wextra -Werror}: four with the GCC and the Clang of Ubuntu 26.04
on x86\_64 and on arm64, and one with the AppleClang of macOS 26 arm64. Four
jobs do so with \option{/W4 /WX}, with the Visual C++ of Visual Studio 2026: x64
in Release, and again asking for the AVX-512 path (compiled where the runner has
the instructions), x64 in Debug, and x86 in
Release. On arm64, the check of the installed headers below compiles the header
of the NEON intervals, which \file{gaol/core/gaol\_interval\_core.h} includes,
with the warnings of the headers (\option{-Wold-style-cast},
\option{-Wundef}) that the intrinsics of \file{<arm\_neon.h>} have to stay clear
of.\newinvfive{} The
first five check the headers GAOL installs with
\file{.github/scripts/headers.sh}: each header compiles alone with
\option{-std=c++11}, the oldest standard GAOL takes, and
\option{-Wall -Wextra -Wold-style-cast -Wundef -Werror}, but those that other
headers include in the middle of their code or for Visual C++ only; and the
macros that \file{gaol/gaol} and the two headers of the parser define, the
include guards apart, start with \code{GAOL\_} or \code{gaol\_}.
\option{-Werror} and \option{/WX} go to the targets only
(\option{CMAKE\_COMPILE\_WARNING\_AS\_ERROR}), so that no check of CMake
fails on a warning of its test program. The library is also compiled with
\option{-Wall -Wconversion} in every build, which fails no job.

\paragraph{The locales and the manual.} The jobs on Ubuntu and in the Debian
containers generate \code{fr\_FR.UTF-8} before the tests with
\cmd{localedef} (\file{.github/scripts/comma-locale.sh}), their systems
having no locale writing a decimal comma; macOS has \code{fr\_FR.UTF-8}, and
Windows \code{French\_France.1252}. After the tests, the jobs of Ubuntu, of the Debian containers, of macOS (but the universal builds and the generator of Xcode), of Visual Studio, of MinGW-w64 and of MSYS2 fail if the output of \code{numbers} holds no check under such a locale,
read in \file{Testing/Temporary/LastTest.log}, \file{tests/numbers.log} with
the autotools or \file{meson-logs/testlog.txt} with meson.\newinvfive{}
Alpine and manylinux have no such locale to generate, and check nothing of
the kind. \code{numbers} has 300~s (\option{GAOL\_NUMBERS\_TIMEOUT}), 400 in
the four Linux jobs with the sanitizers, where it took 109 to 287~s and once more
than 300, and 600 in the job of macOS 15 x86\_64 GCC with the sanitizers,
where it took 146 to 367~s on the same code, and once more than 400. The manual of GAOL v5 is
built with the LaTeX of Ubuntu 24.04, by the autotools, the CMake and the
meson builds, when \file{manual/} changes (\file{manual.yml}), and its PDF,
\file{manual/Manual\_GAOLv5.pdf}, is an artefact of the run; the manual of
GAOL~4 is not built; the examples of this manual that show what they
print are compiled by \file{manual/check\_examples.py} in the job Ubuntu
24.04 x86\_64 GCC at each run (\pxref{sec:examples}).

\subsection{The speed of a pull request}\label{sec:ci-speed}

The times of the continuous integration cannot be compared from one run to
the next. The runners are virtual machines shared with other jobs, whose
speed changes by several times between two runs of the same code:
\code{x * y} with Visual C++ x64 took 7.9~ns and then 16.3, and \code{sqrt}
with MSYS2 and GCC 16 6.8~ns and then 170. The step ``Performance'' of the
jobs prints the table of \code{gaol\_performance} in their summary, and
never fails a job: it tells what a platform costs at that moment, and
nothing else. Showing the best of several passes would not help: a pull
request is to be compared with its base, which is not measured at the same
moment.

\medskip
\file{performance.yml} measures two builds in turn, in the same job, so that
what the runner does between two moments reaches both.\newinvfive{} It runs
for a pull request that changes what is measured, not for one that changes
the documentation: \file{gaol/}, \file{3rd/}, \file{tests/performance.cpp},
the files of the benchmark of \file{doc/compare/code/}, \file{CMakeLists.txt}
and \file{cmake/}, and the scripts and the workflow themselves. A job:
\begin{enumerate}
\item checks out the merge commit of the pull request and its first parent,
  the base it is merged into: the difference between the two is the pull
  request alone, whatever the base branch has gained since the pull request
  was started;
\item configures both with CMake in Release and builds \code{gaol\_performance}
  alone, with the options of the job; the programs of the base are kept
  from one run to the next, keyed by the job, the image of the runner, the
  model of the processor, whether it has the AVX-512 instructions
  (processors of one model do not all have them) and the commit of the
  base, so that the next push of the same pull request builds only the pull
  request, and the two builds come from the same tools and the same
  choices;
\item runs \code{gaol\_performance} three times with each build, the base
  first in the first and third runs and the pull request first in the second
  (\file{.github/scripts/performance-ab.sh});
\item writes in the summary of the job the table of the median of each row
  for the base and for the pull request, their ratio (base over pull
  request: above 1, the pull request is faster) and the spread of the runs,
  the largest ratio of the slowest run to the fastest of either build
  (\file{.github/scripts/compare-performance.py}).
\end{enumerate}

\noindent
The jobs are Ubuntu 24.04 x86\_64 with GCC, with Clang, with GCC and the
rounding direction preserved, and with GCC and with Clang and the FPU
intervals; Ubuntu 24.04 arm64 with GCC and with Clang (the NEON intervals)
and with GCC and the FPU intervals; macOS 26 arm64 and macOS 15 x86\_64;
Visual Studio 2026 x64, x86, arm64, x64 with the direction preserved, and x64
with the direction preserved and the AVX-512 path left out, which measures the
macros that write MXCSR alone, and clang-cl for x64, x86 and arm64; MSYS2 UCRT64 and CLANG64; and Debian 13 containers for i386 and armhf, with no cache. The
other platforms (POWER, s390x, RISC-V, MinGW-w64 of Chocolatey, MSYS2 CLANGARM64
and MINGW32) are not compared, nor are the autotools and meson builds:
the three builds give the same flags, which the audit checks.

\medskip
The table never fails a job because of a time. Only a large difference, on
rows that share their code, means something on a shared runner (the same
code has given a ratio of up to 1.17), and a ratio within the spread of the
runs is noise: under the table, the rows whose ratio is 1.5 or more, or
$2/3$ or less, and beyond the spread of the runs are listed, or none is. A
job fails when a build or a program does, and when the tables cannot be
compared: an output with no row, or no row common to the base and the pull
request. The comparison has its test, which needs Python alone and runs
first in each job (\file{.github/scripts/test-performance-tools.py}): it
fails if the median is replaced by the mean or the threshold is inverted. On
Linux, macOS and with Visual Studio the job also runs the ten long lines of
the benchmark of the comparison with the other libraries
(\pxref{sec:comparison}) with \code{gaol\_bench}, the program \cmd{make perf}
builds too: 200\,000 intervals drawn by \file{bench.py}, 3 repeats of each
line, 3 runs of each build, and a second table under the first; the lines
are left out, with a warning, where one of the builds has no
\code{gaol\_bench} that builds or numpy cannot be installed. The job
\code{control-register-performance} of the containers (Debian 13 armhf and
arm64) uses the same scripts to compare two builds of one tree that differ
by a CMake option: the control register of the rounding direction written
directly with \option{GAOL\_PRESERVE\_ROUNDING}, or the functions of the C
library, on the rows that the register can change. The same comparison runs
anywhere with CMake, a C++ compiler and, for the long lines, Python 3 with
numpy:
\begin{cmdshell}
% sh .github/scripts/performance-ab.sh --before-src ../base --after-src . --lines
% python3 .github/scripts/compare-performance.py base perf-before-1.md \
    perf-before-2.md perf-before-3.md -- PR perf-after-1.md perf-after-2.md \
    perf-after-3.md
\end{cmdshell}
\option{-{}-options} and \option{-{}-before-options} give CMake options to
both builds or to the base alone; the script describes them all.

\subsection{Configurations refused or left out}\label{sec:ci-refused}
\cindex{continuous integration!configurations refused}%

Table~\ref{tab:ci-refused} gives the configurations GAOL v5 refuses, and
those the continuous integration leaves out. The builds refuse the compilers
among them when configuring, or \file{gaol/gaol\_config.h} when compiling
(\pxref{sec:refused}), and the last column says which job checks the
refusal. What a configuration gives with GAOL v5 was measured on 21
September 2026, building it past its refusal.

\begingroup\renewcommand{\arraystretch}{1.1}\small
\begin{longtable}{L{2.9cm}L{5.8cm}L{5cm}}
  \caption{The configurations refused or left out}\label{tab:ci-refused}\\
  \hline
  \multicolumn{1}{c}{Configuration} & \multicolumn{1}{c}{Why} &
  \multicolumn{1}{c}{What GAOL v5 does}\\
  \hline\hline
  \endfirsthead
  \multicolumn{3}{l}{\textit{Table~\thetable, continued}}\\*
  \hline
  \multicolumn{1}{c}{Configuration} & \multicolumn{1}{c}{Why} &
  \multicolumn{1}{c}{What GAOL v5 does}\\
  \hline\hline
  \endhead
  \hline
  \endfoot
  Clang for 32-bit ARM (Debian armhf) & Wrong bounds: Clang does not honour
    the rounding direction there. Built by Clang 19 (Debian 13) and Clang 21
    (Debian sid), GAOL v5 fails \code{arithmetic}, \code{elementary},
    \code{reverse}, \code{other\_functions} and \code{rounding\_direction};
    the products are not the tightest in any direction. & Refused by the
    three builds and by \file{gaol/gaol\_config.h}; a job checks that CMake,
    \cmd{configure} and meson refuse it, for the C compiler alone as for the
    C++ compiler alone. GCC builds the armhf jobs.\\
  Clang 14 on 64-bit ARM, and any compiler saying of
    \option{-frounding-math} ``overriding currently unsupported rounding
    mode on this target'' & Wrong bounds: Clang 14 does not honour the
    rounding direction on 64-bit ARM, and says so. GAOL v5 fails
    \code{arithmetic}, \code{elementary}, \code{reverse} and
    \code{other\_functions} in Release, and \code{elementary} in Debug.
    Clang 18 honours it. & Refused by the three builds; a job checks it for
    the C compiler alone, the C++ compiler alone and both. Ubuntu 22.04
    arm64 is built with GCC only.\\
  MinGW-w64 whose \code{fma()} or \code{round()} is wrong: on x64, GCC 11 to
    13 of Chocolatey (mingw-w64 before 12) and the mingw-w64 linked with
    \file{msvcrt.dll} (MSYS2 MINGW64, the cross compilers of Debian and
    Ubuntu); on 32-bit x86, GCC 11 of WinLibs (mingw-w64 9) and, without
    \option{-mfma}, mingw-w64 12 or later linked with \file{msvcrt.dll}
    (MSYS2 MINGW32); on ARM, as on x64 & Wrong bounds. On x64 the
    \code{fma()} of mingw-w64's own math library adds the products of the
    halves of its arguments with four roundings (10.7\,\% of the error-free
    products and 25 to 73\,\% of other triples wrong), and CORE-MATH
    computes with it wherever the compiler has no \option{-mfma}; its
    \code{round()} depends on the direction. GAOL v5 fails
    \code{elementary}, \code{reverse} and \code{core\_math}: bounds of
    \code{tan()}, \code{asin()}, \code{atan()}, \code{sinh()}, \code{cosh()},
    \code{atanh()} and \code{atan2()} of small arguments do not enclose the
    exact values. From mingw-w64 12 on, a toolchain linking the UCRT takes
    the \code{fma()} and \code{round()} of \file{ucrtbase.dll}. & Refused by
    \file{gaol/gaol\_config.h}; seven jobs check the refusal, three of them
    MSYS2. MinGW-w64 GCC 14 and 15 on x64, 12 to 15 on 32-bit x86 (13 also
    with \option{GAOL\_FMA=OFF}), MSYS2 UCRT64 and CLANG64 are built and
    tested; \code{core\_math} checks \code{fma()} and \code{round()}
    first.\\
  Visual C++, and clang-cl from Clang 16, without \option{/fp:strict}
    (\option{/fp:precise}, their default) & Bounds not certified: they then
    assume rounding to nearest, and may evaluate or rewrite floating-point
    operations accordingly, while GAOL computes rounding upward. & Refused
    by \file{gaol/gaol\_config.h}; \code{gaol::gaol} gives
    \option{/fp:strict} to the code linking it, and each Visual Studio job
    checks the refusal (\file{tests/fp\_strict}), those of clang-cl
    included.\\
  The SSE2 intervals on 32-bit Windows, compiled without aligned \code{new}
    & Crash: GCC takes the memory of \code{new} to be aligned on 16 bytes,
    which the C runtime of 32-bit Windows aligns on 8, and a
    \code{std::vector} of SSE2 intervals crashes on its first
    \code{movaps}; with MinGW-w64 14.2 and 15.2 for x86, \code{arithmetic},
    \code{other\_functions}, \code{reverse} and \code{rounding\_direction}
    crash. Visual C++ and clang-cl have no such flag: only
    \option{/std:c++17} turns the aligned \code{new} on. & The three builds
    give \option{-faligned-new=8} to the library and to the code using it
    (\code{gaol::gaol}, \file{gaol.pc}, \code{gaol\_dep}), and
    \file{gaol/gaol\_config.h} refuses the intervals compiled without
    (\code{refused\_aligned\_new*}). Visual C++ and clang-cl keep the two
    bounds as two doubles of alignment 8
    (\code{GAOL\_SSE2\_UNALIGNED\_STORAGE}), which
    \code{test\_intervals\_\allowbreak aligned\_on\_8\_bytes} exercises.\\
  The NEON intervals on 32-bit ARM & Not offered: they need the vectors of
    two doubles of 64-bit ARM, and write the rounding direction into FPCR
    with \code{mrs} and \code{msr}. & The FPU computes the intervals, as
    before, and \file{gaol/gaol\_config.h} stops the compilation of NEON
    intervals elsewhere; the Visual Studio and MSYS2 CLANGARM64 jobs check
    the NEON intervals on arm64, and the armhf jobs the FPU ones.\\
  Cygwin without the fused multiply-add instructions
    (\option{GAOL\_FMA=OFF}, or a processor without them) & Bounds not
    certified: the \code{fma()} of the C library of Cygwin is not correctly
    rounded, \code{fma(1 + 2\textasciicircum{}-30, 1 - 2\textasciicircum{}-30, -1)} giving 0 instead of
    $-2^{-60}$, and CORE-MATH computes with \code{fma()}. & Refused by
    \file{gaol/gaol\_config.h}; the Cygwin job checks the refusal, and
    builds and tests Cygwin with \option{-mfma} (GCC 14.4, Release, three
    tests skipped).\\
  Doubles computed on the x87 unit of 32-bit x86 & Wrong results rounded to
    nearest, and wrong bounds: the x87 unit computes with 64-bit
    significands, and CORE-MATH assumes every operation on doubles rounded
    to a double. With GCC 12 on Debian 12 i386, CORE-MATH rounded to nearest
    gives the other neighbour of the exact value at 175 of the arguments
    compared with SSE2; with GCC 9 on x86-64 (\option{-mfpmath=387}), the
    bounds of \code{exp2(-1075)}, \code{exp10(-400)}, \code{expm1(-800)} and
    of \code{atan2()} of a tiny and a huge number do not enclose the exact
    values. & Refused by \file{gaol/gaol\_config.h}; the builds compile with
    \option{-msse2 -mfpmath=sse} on 32-bit x86, and
    \code{extended\_precision} fails on the x87 unit.\\
  Debian 11 Bullseye (amd64, arm64, armhf) & Cannot be installed: out of
    support since August 2026, its security repository lists packages it no
    longer serves, so that \cmd{g++} cannot be installed in the container.
    GAOL v5 builds and passes its tests there, with GCC 10.2 and CMake
    3.18.4 installed from snapshot.debian.org. & Left out of the continuous
    integration. Debian 12 and 13 are built.\\
  The intervals of floats (\code{gaol::\allowbreak intervalf},
    \code{gaol::\allowbreak interval2f})
    & Unfinished: \code{sqrt(intervalf)} returns its argument and
    \code{interval2f::inverse()} aborts; neither IBEX nor Codac uses them.
    & Compiled by none of the three builds, which have no option for them:
    only a developer of GAOL defines \code{GAOL\_FLOAT\_INTERVALS}. The
    audit still compares it between the builds.\\
\end{longtable}
\endgroup

\subsection{What these platforms showed}\label{sec:ci-showed}

Building and testing GAOL on these platforms, while the workflows were
written and since their first run on 13 September 2026, showed the
following, which GAOL v5 changed; the register of the differences gives each
change with its measures (Chapter~\ref{chap:changes}).

\begin{itemize}
\item[$\bullet$] \textbf{Bounds that were wrong on some platforms only.}
  The hyperbolic functions, taken from the libm of the system: the libms of
  glibc 2.31, musl and MinGW-w64 are sometimes more than one double from the
  exact value; they are now those of CORE-MATH, correctly rounded, on every
  platform.\newinvfive{} The numbers read from text: the C runtime of
  Windows and musl on 64-bit ARM round \code{strtod()} to nearest in every
  direction, and \code{interval("0.1")} did not enclose $1/10$ there;
  numbers are now read exactly, and intervals written in decimal are
  rounded outward by GAOL rather than by the C library.\newinvfive{}
  \code{mid()} with Visual C++ (\option{/O2 /fp:strict}), which rewrote
  $(-0.5)x+y$ into $y-0.5x$: \code{mid([2\textasciicircum{}-1074, MAX])} was wrong in the
  seven Release jobs of Visual Studio. \code{mid()} on 64-bit ARM with
  autotools and meson, which did not give \option{-ffp-contract=off}: the
  compilers fused its multiplications and additions, and it returned a
  single double; the three builds now give the same flags, which the audit
  jobs check on each kind of machine. Square roots with Visual C++ for
  32-bit x86, whose \code{sqrt} rounds to nearest in every rounding
  direction.
\item[$\bullet$] \textbf{Operations that were slow on some platforms only},
  which the table of \code{gaol\_performance} in the summary of each
  Release job showed. Reading the rounding direction before each operation
  made \code{x + y} take 800~ns under Rosetta 2 and 82 to 98~ns with 32-bit
  Visual C++: an addition shows whether it is upward instead.
  \code{fesetround()} cost 130~ns with mingw-w64 13, and 50 to 250~ns with
  Visual C++: on x86 processors the control registers are written directly.
  \option{-ffloat-store}, given to GCC on every target, made \code{x * y}
  take 18.3~ns rather than 4.9 on x86-64: it is given only where doubles are
  still computed on the x87 unit.
\item[$\bullet$] \textbf{Crashes and builds that failed.} \cmd{configure}
  and meson did not install the headers for MinGW and Visual C++, and meson
  did not compile on Linux. The agreement jobs passed while a build had not
  configured in the containers, a pipe hiding the failure: they now fail,
  and print the logs. \cmd{configure -{}-enable-debug} did not compile, the
  Debug jobs of CMake now define \code{GAOL\_DEBUGGING} too, and the test
  \code{debugging} compiles that code. \code{minimum()} and \code{maximum()}
  of two zeros of different signs gave $+0$ and $-0$ with Visual C++ in
  Release, rather than $-0$ and $+0$, which the unit test
  \code{float\_functions} checks, built with Visual C++ since the tests of
  \file{check/} are in \file{tests/}. A shared \file{libgaol} built by
  \cmd{configure} or meson did not export the classes and functions of the
  expressions, and a program building one did not link with it: the shared
  library jobs of the three builds, whose tests link with the installed
  \file{libgaol.so}, build \file{tests/expressions.cpp}.
\item[$\bullet$] \textbf{Reports of the sanitizers that were suppressed at
  first}, and are fixed now: the nodes the parser did not free. Nothing is
  suppressed any more.
\end{itemize}

\section{The scripts of the repository}\label{sec:scripts}
\cindex{scripts}%

Every script of the repository starts with a comment that says how to run it
and what it needs. Tables \ref{tab:scripts-generated} to
\ref{tab:scripts-tools} give what each one is for, so that the right one is
found without opening them all; the scripts of the comparisons with the
other libraries (\file{doc/compare/code/}, \file{doc/compare/itf1788/},
\file{doc/compare/tang2021/}) are described with the comparisons
(\pxref{sec:comparison} and \exref{sec:test-suites}). The files taken from
other projects, which GAOL does not change nor document one by one, are the
scripts of the autotools at the root (\file{config.guess},
\file{config.sub}, \file{configure}, \file{depcomp}, \file{compile},
\file{install-sh}, \file{ltmain.sh}, \file{missing}, \file{test-driver})
and the tools and the CI of CORE-MATH under \file{3rd/math-core/}.

\begingroup\renewcommand{\arraystretch}{1.1}\small
\begin{longtable}{L{5cm}L{8.8cm}}
  \caption{The scripts of the generated files, the manual and the coverage}
  \label{tab:scripts-generated}\\
  \hline
  \multicolumn{1}{c}{Script} & \multicolumn{1}{c}{What it is for}\\
  \hline\hline
  \endfirsthead
  \multicolumn{2}{l}{\textit{Table~\thetable, continued}}\\*
  \hline
  \multicolumn{1}{c}{Script} & \multicolumn{1}{c}{What it is for}\\
  \hline\hline
  \endhead
  \hline
  \endfoot
  \file{gaol/parser/regenerate\_parser.sh} & Regenerates the lexer and the
    parser from the \file{.lpp} and \file{.ypp} sources (flex 2.6.4, bison
    3.8.2), after any change of them (\cmd{make parser})\\
  \file{gaol/float/gaol\_generator\_mul\_i2f.py} & Writes the 256 cases of the
    multiplication of \code{interval2f} by the signs of the operands: a tool for the
    developers (Python~3 and \code{ply}), which no build runs\\
  \file{manual/build-pdf.sh} & Builds the manual v5 with \cmd{pdflatex}, in
    \file{manual/Manual\_GAOLv5.pdf}; \cmd{make pdf} and the target
    \code{pdf} of CMake and meson call it\\
  \file{manual/check\_examples.py} & Compiles and runs each example of the
    manual v5 with an installed GAOL, and compares what it prints with what
    the manual shows\\
  \file{coverage/coverage\_summary.py} & Reads the JSON of gcovr and writes
    \file{coverage/README.md}: lines and branches by file, GAOL apart from
    CORE-MATH, and the lines never run\\
  \file{cmake/gaol\_version.cmake} & \code{gaol\_read\_version()} reads
    \file{VERSION.txt} before \code{project()}, as \file{configure.ac} and
    \file{meson.build} do\\
  \file{cmake/gaol\_package\_source.cmake} & Included by CPack: brings the
    name and the version of the archive of the sources up to date after a
    change of \file{VERSION.txt}\\
\end{longtable}
\endgroup

\begingroup\renewcommand{\arraystretch}{1.1}\small
\begin{longtable}{L{5cm}L{8.8cm}}
  \caption[The audit of the three builds and the checks of the continuous \mbox{integration}]{The audit of the three builds (\file{.github/audit/}) and the
    checks of the continuous integration (\file{.github/scripts/})}
  \label{tab:scripts-ci}\\
  \hline
  \multicolumn{1}{c}{Script} & \multicolumn{1}{c}{What it is for}\\
  \hline\hline
  \endfirsthead
  \multicolumn{2}{l}{\textit{Table~\thetable, continued}}\\*
  \hline
  \multicolumn{1}{c}{Script} & \multicolumn{1}{c}{What it is for}\\
  \hline\hline
  \endhead
  \hline
  \endfoot
  \file{audit/audit.sh} & Configures GAOL with CMake, autotools and meson,
    and writes a JSON file of what each sees (with \file{probe\_build.py})\\
  \file{audit/probe\_build.py}, \file{audit/make\_probe.py} & Read the
    compiler flags and the configuration GAOL sees in the three builds, for
    a C++ source and for a C source of CORE-MATH, through a file that shows
    the configuration once preprocessed\\
  \file{audit/compare.py} & Compares the macros and the flags of the three
    builds, and with \option{-{}-check} fails if they differ\\
  \file{audit/logs.sh} & When the audit fails, prints the end of the logs of
    the three builds, then exits with the same status\\
  \hline
  \file{scripts/tests.sh} & Builds the tests with a GAOL installed by
    autotools or meson (static or shared), then runs them\\
  \file{scripts/headers.sh} & Checks that each installed header compiles
    alone in C++11 with
    \option{-Wall -Wextra -Wold-style-cast -Wundef -Werror}\\
  \file{scripts/\allowbreak check-intervals.sh} & Fails unless the build computes with
    the expected intervals (NEON, SSE2 or FPU)\\
  \file{scripts/\allowbreak check-gcc-windows-fma.sh} & Checks that the objects of \file{gaol\_interval\_core.cpp} and of the \code{cos} of CORE-MATH built by GCC for Windows hold fused multiply-add instructions, and that no object names a 256-bit or 512-bit register (an error on x64, GCC bug 54412; a notice on x86)\\
  \file{scripts/comma-locale.sh} & Generates \code{fr\_FR.UTF-8} with \cmd{localedef} before the tests (\code{generate}), and checks after them that \code{numbers} ran checks under a locale writing a
decimal comma (\code{check}), so that \code{numbers} does not pass without having tried a locale writing a decimal comma\\
  \file{scripts/sde-tests.sh} & Runs the tests that compute with intervals (all
    but \code{numbers}, \code{core\_math} and \code{u128}) under Intel SDE,
    which emulates AVX-512: the runners of GitHub for Windows may have none\\
  \file{scripts/version-file.sh} & Checks how \cmd{configure}, meson and
    autoconf read \file{VERSION.txt} written in several ways\\
  \file{scripts/performance-ab.sh} & Builds and measures two builds in turn
    in the same job: the base and the head of a pull request, or one source
    with two options\\
  \file{scripts/\allowbreak compare-performance.py} & Compares two series of measures
    (median, ratio, spread) and writes a Markdown table for the summary of
    a job\\
  \file{scripts/\allowbreak test-performance-tools.py} & Tests
    \file{compare-performance.py} and the name of the processor, in Python
    alone and without a build\\
\end{longtable}
\endgroup

\begingroup\renewcommand{\arraystretch}{1.1}\small
\begin{longtable}{L{5cm}L{8.8cm}}
  \caption{The scripts of the tests (\file{tests/}) and their tools
    (\file{tests/tools/})}\label{tab:scripts-tools}\\
  \hline
  \multicolumn{1}{c}{Script} & \multicolumn{1}{c}{What it is for}\\
  \hline\hline
  \endfirsthead
  \multicolumn{2}{l}{\textit{Table~\thetable, continued}}\\*
  \hline
  \multicolumn{1}{c}{Script} & \multicolumn{1}{c}{What it is for}\\
  \hline\hline
  \endhead
  \hline
  \endfoot
  \file{elementary\_values.py} & Writes \file{elementary\_values.h}: for
    each $f(x)$, the greatest double below and the least double above,
    computed with mpmath, for \code{elementary}\\
  \file{extended\_precision\_values.py} & Writes
    \file{extended\_precision\_values.h}: the bounding values and the
    nearest double, to compare CORE-MATH and GAOL\\
  \file{compoundm1\_values.py} & Writes \file{compoundm1\_values.h}: the
    doubles around $(1+x)^y-1$ at pairs, the hull over boxes, and
    $\log(1+x)$ and $e^t-1$ to 256 bits, with mpmath and fractions\\
  \file{reverse\_values.py}, \file{atan2\_rel\_values.py},
    \file{pow\_rel\_values.py} & Write the cases of the reverse functions of
    IEEE 1788-2015, of \ieee{atan2Rev1} and \ieee{atan2Rev2}, and of
    \ieee{powRev1} and \ieee{powRev2}, for \code{reverse}\\
  \file{refused\_options.sh} & Checks that autotools and meson refuse
    \option{-ffast-math} and \option{-ffinite-math-only}, and the SSE2
    intervals of 32-bit Windows without \option{-faligned-new=8}, as the tests
    \code{refused\_*} and \code{refused\_aligned\_new*} of CMake do\\
  \file{release\_flags.py} & Checks that a project which brings GAOL in
    without a build type compiles it in release (\option{-O3},
    \code{NDEBUG}, \code{GAOL\_DEBUGGING} at 0)\\
  \file{meson\_subproject/check\_flags.py} & Checks that a meson project
    which takes GAOL as a subproject gets the flags of its \file{gaol.pc}\\
  \file{compile\_error.cmake} & A \cmd{cmake -P} script that compiles a
    source whose compilation has to fail (\code{refused\_*},
    \code{nodiscard\_discard\_cxx*})\\
  \file{cpack\_stale\_configure.cmake}, \file{automatic\_configuration.cmake},
    \file{version\_file.cmake} & The three tests of the CMake build that are
    scripts (\pxref{sec:test-programs})\\
  \hline
  \file{tools/arith/diffarith.py} & Differential test of the arithmetic:
    the bounds and the exception flags of \code{+ - * / \%}, \code{sqr},
    \code{sqrt}, \code{inverse}, \code{div\_rel} and \code{mid} on a grid of
    boxes, under the four rounding directions, with two builds of GAOL\\
  \file{tools/pow/diffpow.py} & Differential test of \code{pow}: the bounds
    of \code{gaol\_ieee1788::pow} and \code{gaol::pow} with two builds, on
    70\,500 boxes and under the four directions\\
  \file{tools/pow/build.py} & Compiles a program of \file{tests/tools/pow}
    against a build of GAOL, with the flags that build gives its own tests\\
  \file{tools/pow/negative\_power\_rows.py} & Writes the rows of the table of
    large negative powers of \file{tests/arithmetic.cpp}: the tightest bounds
    of $x^n$ for $n$ from $-2^{31}+1$ to $-2^{27}$, computed with mpmath on
    2000 bits\\
  \file{tools/pow/gen\_table.py}, \file{tools/pow/checkrows.py} & Read and
    rewrite the table of boxes of \code{ieee1788}, and check its bounds
    against the exact power (mpmath, 500 bits)\\
  \file{tools/pow/mutate.py}, \file{tools/cos/mutate.py},
    \file{tools/x86-csr/mutants.py} & Mutation testing: one deliberate error
    in \code{pow} and the integer powers, in the port of \file{cos.c} to the
    128-bit integer, or in the writes of the control registers of x86, and
    the tests have to fail\\
  \file{tools/pow/pown\_bound.py} & Checks with rationals the bound of the
    error of the integer powers computed beforehand\\
  \file{tools/pow/bench.cpp}, \file{tools/pow/pownhash.cpp},
    \file{tools/pow/pownedge.cpp} & The time per call of each path of
    \code{pow}, and a hash of the bounds of four million integer powers, the
    sign of a zero included\\
  \file{tools/cos/accurate\_values.py} & Computes the roundings of
    $\cos x$ with mpmath and checks the table of the accurate path of
    \code{core\_math}\\
  \file{tools/asm-order/check\_order.py} & Reads the assembly of the library
    (x86-64, GCC and Clang) for an instruction that rounds on the wrong side
    of a write of the control register (GCC bug 34678), and for the
    conversion of a subnormal constant through the x87 unit\\
  \file{tools/ibex\_codac/*.py} & The cases of the tables of
    \code{test\_ibex\_codac}, their exact results (rationals, mpmath), and
    the check that the rows of the test are those the cases give\\
  \file{tools/x86-csr/\allowbreak amalgamate.py} & One source file from three units of
    GAOL, to compile with a compiler that takes one file, such as those of
    Compiler Explorer\\
\end{longtable}
\endgroup

\noindent
The workflows are read by actionlint, which runs \cmd{shellcheck} on the
steps that have a script, and reports nothing; the scripts of
\file{.github/} and of \file{doc/compare/code/} are read by
\cmd{shellcheck}, run from the root of the repository, which reports
nothing either: the lists of flags that the shell has to split carry a
comment with the reason.\newinvfive

\section{Coverage}\label{sec:coverage}
\cindex{coverage}%

The option \option{GAOL\_COVERAGE} of CMake (\loption{enable-coverage} of
\cmd{configure}, \code{enable-coverage} of meson), off by default, compiles
GAOL and its tests with the counters of gcov, and adds the target that runs
the tests and writes the report of their coverage,
\file{coverage/coverage.html} and the page of conclusions
\file{coverage/README.md}: \cmd{cmake -{}-build <build> -{}-target coverage},
\cmd{make coverage}, or \cmd{ninja gaol-coverage}, meson keeping
\code{coverage} for its own report (\pxref{sec:cmake-options}).\newinvfive{}
It needs GCC or Clang, gcovr and Python~3, and the tests (\option{WITH\_TESTS}); it should be the Debug build: with the Release build, the optimizations move the code around, the counts no longer follow the source, and the build warns. The sources of GAOL
(\file{gaol/}) are counted apart from those of CORE-MATH
(\file{3rd/math-core/}), whose accurate phases are reached by a handful of
arguments in a million: counting them together would say nothing of either.
Table~\ref{tab:coverage} gives the report of the run of 4 November 2026.

\begin{table}[htbp]
  \centering\renewcommand{\arraystretch}{1.1}
  \caption{The coverage of the tests on 4 November 2026} \label{tab:coverage}
  \begin{tabular}{lrrrr}
    \hline
    & Lines & Covered & Not run & \%\\
    \hline\hline
    GAOL (\file{gaol/}) & 5466 & 5153 & 313 & 94.3\\
    CORE-MATH (\file{3rd/math-core/}) & 8172 & 6732 & 1440 & 82.4\\
    \hline
  \end{tabular}

\end{table}

\noindent
The tests run 94.3\,\% of the lines of GAOL, above the 80\,\% GAOL v5 aims
at, and 66.5\,\% of its branches are taken both ways. The lines of CORE-MATH
that are never run are its accurate phases, which a handful of arguments in
a million reach, and the paths of the architectures the machine is not: they
are counted among the lines not run, and are checked otherwise: the accurate
path of \code{sin} is compared with the upstream sources compiled unchanged,
and that of \code{cos} with its roundings computed by mpmath
(\pxref{sec:test-programs}), and the other accurate phases are reached by the
tests of their bounds only. File by file, the least
covered are \file{gaol/parser/gaol\_expr\_visitor.h} (9 of 99 lines, 9.1\,\%),
\file{gaol/parser/gaol\_expr\_eval.h} (87.3\,\%),
\file{gaol/parser/gaol\_expression.cpp} (89.6\,\%) and
\file{gaol/core/gaol\_interval\_sse.cpp} (93.9\,\%);
\file{gaol/core/gaol\_interval\_core.cpp}, the largest file, with 2492 lines counted,
is covered at 97.0\,\%, and 18 files of the 26 at 100\,\%. The page lists
the lines the tests never run, file by file.

\section{Discussion}

The tests of GAOL v5 compare its bounds with references that owe nothing to
the library nor to the floating-point environment, and run every operation
in every state a program may leave the processor in: the four rounding
directions, the modes that flush the subnormal numbers, the exceptions
enabled, a locale with a decimal comma, and the initialization of static
objects before \code{main()}. The programs of \file{tests/} hold hundreds
of thousands of such checks, and the documentation records, for many of the changes, how many of the checks written with them fail without them, and the example
programs check what they print, as do the examples of this manual. The continuous integration runs
the suite on three operating systems, seven processor architectures and the
compilers of each, with the three builds, and refuses the configurations
whose bounds it found wrong or uncertified, most of the refusals being checked by a
job of its own.

\medskip
Three choices of the tests and of the continuous integration deserve a word. The teeth of some checks depend
on the processor the code is compiled for: a loop that a compiler
vectorizes raises the invalid-operation flag with one set of instructions
and not with another, so that the jobs on many processors are the test, not
one of them. The speed of a pull request cannot be read from the times of a
shared runner, which change by several times between two runs of the same
code: it is measured by building the base and the head in the same job and
running them in turn, and only a ratio of 1.5 or more, or of 2/3 or less, beyond the spread of
the runs, means something. The coverage report counts the lines of GAOL
apart from those of CORE-MATH, whose accurate phases a handful of arguments
in a million reach; \code{core\_math} checks the accurate paths of \code{sin}
and \code{cos} against values computed apart instead.

\medskip
Some checks are deferred or left out, for the reasons the documentation
gives: the reading of a text of 512~MiB, which found a 32-bit exponent
overflow, has no test, since the text is read through two more copies of
it; the checks of several threads reading or writing at once are commented
out, the tests running no thread; the refusal of the mingw-w64 of ARM whose
\code{fma()} is wrong is not tested; Debian 11 is left out of the
continuous integration, and the intervals of floats are compiled by none of
the builds. The next chapter turns to the programs written for GAOL~4: how
to compile and link them with GAOL v5, which results change, and the
register of the differences by theme.
